
\documentclass{article}
\usepackage{colortbl}
\usepackage{makecell}
\usepackage{multirow}
\usepackage{supertabular}

\begin{document}

\newcounter{utterance}

\centering \large Interaction Transcript for game `clean\_up', experiment `2\_hard\_7obj\_pl', episode 0 with glm{-}5.2.
\vspace{24pt}

{ \footnotesize  \setcounter{utterance}{1}
\setlength{\tabcolsep}{0pt}
\begin{supertabular}{c@{$\;$}|p{.15\linewidth}@{}p{.15\linewidth}p{.15\linewidth}p{.15\linewidth}p{.15\linewidth}p{.15\linewidth}}
    \# & $\;$A & \multicolumn{4}{c}{Game Master} & $\;\:$B\\
    \hline

    \theutterance \stepcounter{utterance}  
    & & \multicolumn{4}{p{0.6\linewidth}}{
        \cellcolor[rgb]{0.9,0.9,0.9}{
            \makecell[{{p{\linewidth}}}]{
                \texttt{\tiny{[P1$\langle$GM]}}
                \texttt{Jestem mistrzem gry, a Państwo grają w grę kooperacyjną z następującą siatką jako planszą:} \\
\\ 
\texttt{```} \\
\texttt{ 1234567} \\
\texttt{╔══╤════╗} \\
\texttt{║◌◌│◌◌◌A║ 1} \\
\texttt{╟──┤◌◌┌─╢ 2} \\
\texttt{║◌◌├──┤T║ 3} \\
\texttt{╟──┘M◌├─╢ 4} \\
\texttt{║P◌OD◌│S║ 5} \\
\texttt{╟─┐◌┌─┼─╢ 6} \\
\texttt{║◌│◌│◌│◌║ 7} \\
\texttt{╚═╧═╧═╧═╝} \\
\\ 
\texttt{```} \\
\\ 
\texttt{* Górna krawędź wyświetla współrzędne x rosnące w prawo, a prawa krawędź współrzędne y rosnące w dół.} \\
\texttt{* Następujące obiekty są losowo umieszczone na siatce: 'P', 'O', 'T', 'S', 'D', 'A', 'M'.} \\
\\ 
\\ 
\texttt{Drugi gracz widzi wariant planszy, na której obiekty są umieszczone w różnych losowych miejscach. Nie można zobaczyć planszy drugiego gracza, a on nie może zobaczyć Państwa planszy.} \\
\\ 
\\ 
\texttt{**Cel:**} \\
\\ 
\texttt{Obaj gracze muszą przesuwać obiekty na swoim tle tak, aby identyczne obiekty znalazły się na tych samych współrzędnych. Muszą Państwo komunikować się z drugim graczem, aby uzgodnić wspólną konfigurację docelową.} \\
\\ 
\\ 
\texttt{**Zasady:**} \\
\\ 
\texttt{* W każdej turze można wysłać dokładnie jedno z dwóch poniższych poleceń:} \\
\texttt{ 1. `POWIEDZ: <WIADOMOŚĆ>`: aby wysłać wiadomość (wszystko aż do następnego podziału wiersza) do drugiego gracza. Przekażę ją Państwa partnerowi.} \\
\texttt{2. `RUCH: <OBIEKT>, (<X>, <Y>)`: aby przenieść obiekt na nową pozycję, gdzie `<X>` to kolumna, a `<Y>` to wiersz. Poinformuję Państwa, czy przeniesienie się powiodło, czy nie.} \\
\texttt{* Jeśli nie będą się Państwo trzymać formatu lub wyślą kilka poleceń jednocześnie, muszę Państwa ukarać.} \\
\texttt{* Jeśli obaj gracze zgromadzą więcej niż 16 kar, obaj przegrają grę.} \\
\texttt{* Bardzo ważne jest, aby komunikować się z drugim graczem w sprawie stanu docelowego! *Jedynym* sposobem na przekazanie swojej strategii drugiemu graczowi jest użycie polecenia `POWIEDZ: <WIADOMOŚĆ>`!} \\
\\ 
\\ 
\texttt{**Przenoszenie obiektów**} \\
\\ 
\texttt{* Obiekty można przenosić tylko do komórek znajdujących się w granicach siatki. Komórka docelowa musi być pusta, tzn. musi zawierać tylko symbol '◌'.} \\
\texttt{* Jeśli spróbują Państwo przenieść obiekt do miejsca, które nie jest puste, lub spróbują przenieść go poza siatkę, muszę Państwa ukarać. Otrzymają Państwo kolejną próbę.} \\
\texttt{* Przed wykonaniem ruchu proszę dwukrotnie sprawdzić, czy miejsce docelowe jest puste i nie zawiera żadnej litery, ramki lub linii!} \\
\\ 
\\ 
\texttt{**Koniec gry**} \\
\\ 
\texttt{Jeśli uważają Państwo, że osiągnęli cel wyrównania wszystkich obiektów, mogą Państwo poprosić drugiego gracza o zakończenie gry, wysyłając `POWIEDZ: zakończono?`. Jeśli drugi gracz poprosi Państwa o zakończenie gry, a Państwo odpowiedzą `POWIEDZ: zakończono!`, gra się zakończy.} \\
\\ 
\texttt{Obaj gracze wygrywają, jeśli gra zakończy się w ciągu 28 rund, gdzie jedna runda jest definiowana jako wysłanie ważnego polecenia przez dwóch graczy.} \\
\\ 
\texttt{**Punktacja:**} \\
\\ 
\texttt{Im bliżej siebie znajdują się identyczne obiekty na obu planszach, tym więcej punktów otrzymują Państwo. Kary zmniejszają liczbę punktów. Czy uda się Państwu pobić rekord?} \\
\\ 
\\ 
\texttt{Proszę wysłać wiadomość do drugiego gracza, aby rozpocząć grę!} \\
            }
        }
    }
    & \\ \\

    \theutterance \stepcounter{utterance}  
    & \multicolumn{4}{p{0.6\linewidth}}{
        \cellcolor[rgb]{0.9,0.9,0.9}{
            \makecell[{{p{\linewidth}}}]{
                \texttt{\tiny{[P1$\rangle$GM]}}
                \texttt{POWIEDZ: Cześć! Zaproponuj układ docelowy, w którym każdy obiekt znajdzie się w osobnym, pustym polu na krawędzi planszy. Przykładowo: A=(4,1), T=(5,1), M=(2,1), S=(7,1), O=(7,2), P=(7,6), D=(7,7). Dostosuję się do Twojej wersji. Czekam na potwierdzenie.} \\
            }
        }
    }
    & & \\ \\

    \theutterance \stepcounter{utterance}  
    & & & \multicolumn{2}{p{0.3\linewidth}}{
        \cellcolor[rgb]{0.9,0.9,0.9}{
            \makecell[{{p{\linewidth}}}]{
                \texttt{\tiny{[GM$|$GM]}}
                \texttt{POWIEDZ: Cześć! Zaproponuj układ docelowy, w którym każdy obiekt znajdzie się w osobnym, pustym polu na krawędzi planszy. Przykładowo: A=(4,1), T=(5,1), M=(2,1), S=(7,1), O=(7,2), P=(7,6), D=(7,7). Dostosuję się do Twojej wersji. Czekam na potwierdzenie.} \\
            }
        }
    }
    & & \\ \\

    \theutterance \stepcounter{utterance}  
    & & \multicolumn{4}{p{0.6\linewidth}}{
        \cellcolor[rgb]{0.9,0.9,0.9}{
            \makecell[{{p{\linewidth}}}]{
                \texttt{\tiny{[GM$\rangle$P2]}}
                \texttt{Jestem mistrzem gry, a Państwo grają w grę kooperacyjną z następującą siatką jako planszą:} \\
\\ 
\texttt{```} \\
\texttt{ 1234567} \\
\texttt{╔══╤════╗} \\
\texttt{║T◌│D◌◌A║ 1} \\
\texttt{╟──┤◌◌┌─╢ 2} \\
\texttt{║◌◌├──┤◌║ 3} \\
\texttt{╟──┘◌◌├─╢ 4} \\
\texttt{║OP◌◌◌│S║ 5} \\
\texttt{╟─┐◌┌─┼─╢ 6} \\
\texttt{║◌│◌│M│◌║ 7} \\
\texttt{╚═╧═╧═╧═╝} \\
\\ 
\texttt{```} \\
\\ 
\texttt{* Górna krawędź wyświetla współrzędne x rosnące w prawo, a prawa krawędź współrzędne y rosnące w dół.} \\
\texttt{* Następujące obiekty są losowo umieszczone na siatce: 'P', 'O', 'T', 'S', 'D', 'A', 'M'.} \\
\\ 
\\ 
\texttt{Drugi gracz widzi wariant planszy, na której obiekty są umieszczone w różnych losowych miejscach. Nie można zobaczyć planszy drugiego gracza, a on nie może zobaczyć Państwa planszy.} \\
\\ 
\\ 
\texttt{**Cel:**} \\
\\ 
\texttt{Obaj gracze muszą przesuwać obiekty na swoim tle tak, aby identyczne obiekty znalazły się na tych samych współrzędnych. Muszą Państwo komunikować się z drugim graczem, aby uzgodnić wspólną konfigurację docelową.} \\
\\ 
\\ 
\texttt{**Zasady:**} \\
\\ 
\texttt{* W każdej turze można wysłać dokładnie jedno z dwóch poniższych poleceń:} \\
\texttt{ 1. `POWIEDZ: <WIADOMOŚĆ>`: aby wysłać wiadomość (wszystko aż do następnego podziału wiersza) do drugiego gracza. Przekażę ją Państwa partnerowi.} \\
\texttt{2. `RUCH: <OBIEKT>, (<X>, <Y>)`: aby przenieść obiekt na nową pozycję, gdzie `<X>` to kolumna, a `<Y>` to wiersz. Poinformuję Państwa, czy przeniesienie się powiodło, czy nie.} \\
\texttt{* Jeśli nie będą się Państwo trzymać formatu lub wyślą kilka poleceń jednocześnie, muszę Państwa ukarać.} \\
\texttt{* Jeśli obaj gracze zgromadzą więcej niż 16 kar, obaj przegrają grę.} \\
\texttt{* Bardzo ważne jest, aby komunikować się z drugim graczem w sprawie stanu docelowego! *Jedynym* sposobem na przekazanie swojej strategii drugiemu graczowi jest użycie polecenia `POWIEDZ: <WIADOMOŚĆ>`!} \\
\\ 
\\ 
\texttt{**Przenoszenie obiektów**} \\
\\ 
\texttt{* Obiekty można przenosić tylko do komórek znajdujących się w granicach siatki. Komórka docelowa musi być pusta, tzn. musi zawierać tylko symbol '◌'.} \\
\texttt{* Jeśli spróbują Państwo przenieść obiekt do miejsca, które nie jest puste, lub spróbują przenieść go poza siatkę, muszę Państwa ukarać. Otrzymają Państwo kolejną próbę.} \\
\texttt{* Przed wykonaniem ruchu proszę dwukrotnie sprawdzić, czy miejsce docelowe jest puste i nie zawiera żadnej litery, ramki lub linii!} \\
\\ 
\\ 
\texttt{**Koniec gry**} \\
\\ 
\texttt{Jeśli uważają Państwo, że osiągnęli cel wyrównania wszystkich obiektów, mogą Państwo poprosić drugiego gracza o zakończenie gry, wysyłając `POWIEDZ: zakończono?`. Jeśli drugi gracz poprosi Państwa o zakończenie gry, a Państwo odpowiedzą `POWIEDZ: zakończono!`, gra się zakończy.} \\
\\ 
\texttt{Obaj gracze wygrywają, jeśli gra zakończy się w ciągu 28 rund, gdzie jedna runda jest definiowana jako wysłanie ważnego polecenia przez dwóch graczy.} \\
\\ 
\texttt{**Punktacja:**} \\
\\ 
\texttt{Im bliżej siebie znajdują się identyczne obiekty na obu planszach, tym więcej punktów otrzymują Państwo. Kary zmniejszają liczbę punktów. Czy uda się Państwu pobić rekord?} \\
\\ 
\texttt{Drugi gracz rozpoczął grę wysyłając następującą wiadomość:} \\
\\ 
\texttt{"Cześć! Zaproponuj układ docelowy, w którym każdy obiekt znajdzie się w osobnym, pustym polu na krawędzi planszy. Przykładowo: A=(4,1), T=(5,1), M=(2,1), S=(7,1), O=(7,2), P=(7,6), D=(7,7). Dostosuję się do Twojej wersji. Czekam na potwierdzenie."} \\
\\ 
\texttt{Jakie jest twoje pierwsze polecenie?} \\
            }
        }
    }
    & \\ \\

    \theutterance \stepcounter{utterance}  
    & & & \multicolumn{4}{p{0.6\linewidth}}{
        \cellcolor[rgb]{0.9,0.9,0.9}{
            \makecell[{{p{\linewidth}}}]{
                \texttt{\tiny{[GM$\langle$P2]}}
                \texttt{POWIEDZ: Proponuję docelowo: A=(6,1), T=(3,1), M=(5,1), S=(7,1), O=(5,7), P=(6,7), D=(7,2). Pasuje?} \\
            }
        }
    }
     \\ \\

    \theutterance \stepcounter{utterance}  
    & & & \multicolumn{2}{p{0.3\linewidth}}{
        \cellcolor[rgb]{0.9,0.9,0.9}{
            \makecell[{{p{\linewidth}}}]{
                \texttt{\tiny{[GM$|$GM]}}
                \texttt{POWIEDZ: Proponuję docelowo: A=(6,1), T=(3,1), M=(5,1), S=(7,1), O=(5,7), P=(6,7), D=(7,2). Pasuje?} \\
            }
        }
    }
    & & \\ \\

    \theutterance \stepcounter{utterance}  
    & & \multicolumn{4}{p{0.6\linewidth}}{
        \cellcolor[rgb]{0.9,0.9,0.9}{
            \makecell[{{p{\linewidth}}}]{
                \texttt{\tiny{[P1$\langle$GM]}}
                \texttt{Twoja wiadomość została przekazana drugiemu graczowi.} \\
\\ 
\\ 
\texttt{Obecnie grasz rundę 1 z maksymalnie 28.} \\
\texttt{Łącznie zgromadziliście 0 z 16 kar.} \\
\texttt{Drugi gracz wysłał tę wiadomość:} \\
\\ 
\\ 
\texttt{"Proponuję docelowo: A=(6,1), T=(3,1), M=(5,1), S=(7,1), O=(5,7), P=(6,7), D=(7,2). Pasuje?"} \\
\\ 
\\ 
\texttt{Jakie jest twoje następne polecenie?} \\
            }
        }
    }
    & \\ \\

    \theutterance \stepcounter{utterance}  
    & \multicolumn{4}{p{0.6\linewidth}}{
        \cellcolor[rgb]{0.9,0.9,0.9}{
            \makecell[{{p{\linewidth}}}]{
                \texttt{\tiny{[P1$\rangle$GM]}}
                \texttt{POWIEDZ: Pasuje. Zgodziłem się na Twój układ docelowy. Możemy działać niezależnie i nie musimy na siebie czekać. Ja zacznę od obniżenia P z (1,5) do (1,7). Po wszystkim poproszę o zakończenie.} \\
            }
        }
    }
    & & \\ \\

    \theutterance \stepcounter{utterance}  
    & & & \multicolumn{2}{p{0.3\linewidth}}{
        \cellcolor[rgb]{0.9,0.9,0.9}{
            \makecell[{{p{\linewidth}}}]{
                \texttt{\tiny{[GM$|$GM]}}
                \texttt{POWIEDZ: Pasuje. Zgodziłem się na Twój układ docelowy. Możemy działać niezależnie i nie musimy na siebie czekać. Ja zacznę od obniżenia P z (1,5) do (1,7). Po wszystkim poproszę o zakończenie.} \\
            }
        }
    }
    & & \\ \\

    \theutterance \stepcounter{utterance}  
    & & \multicolumn{4}{p{0.6\linewidth}}{
        \cellcolor[rgb]{0.9,0.9,0.9}{
            \makecell[{{p{\linewidth}}}]{
                \texttt{\tiny{[GM$\rangle$P2]}}
                \texttt{Twoja wiadomość została przekazana drugiemu graczowi.} \\
\\ 
\\ 
\texttt{Obecnie grasz rundę 2 z maksymalnie 28.} \\
\texttt{Łącznie zgromadziliście 0 z 16 kar.} \\
\texttt{Drugi gracz wysłał tę wiadomość:} \\
\\ 
\\ 
\texttt{"Pasuje. Zgodziłem się na Twój układ docelowy. Możemy działać niezależnie i nie musimy na siebie czekać. Ja zacznę od obniżenia P z (1,5) do (1,7). Po wszystkim poproszę o zakończenie."} \\
\\ 
\\ 
\texttt{Jakie jest twoje następne polecenie?} \\
            }
        }
    }
    & \\ \\

    \theutterance \stepcounter{utterance}  
    & & & \multicolumn{4}{p{0.6\linewidth}}{
        \cellcolor[rgb]{0.9,0.9,0.9}{
            \makecell[{{p{\linewidth}}}]{
                \texttt{\tiny{[GM$\langle$P2]}}
                \texttt{RUCH: O, (5,7)} \\
            }
        }
    }
     \\ \\

    \theutterance \stepcounter{utterance}  
    & & & \multicolumn{2}{p{0.3\linewidth}}{
        \cellcolor[rgb]{0.9,0.9,0.9}{
            \makecell[{{p{\linewidth}}}]{
                \texttt{\tiny{[GM$|$GM]}}
                \texttt{RUCH: O, (5,7)} \\
            }
        }
    }
    & & \\ \\

    \theutterance \stepcounter{utterance}  
    & & & \multicolumn{2}{p{0.3\linewidth}}{
        \cellcolor[rgb]{0.9,0.9,0.9}{
            \makecell[{{p{\linewidth}}}]{
                \texttt{\tiny{[GM$|$GM]}}
                \texttt{Kara: (5,7) nie jest puste, ale zawiera 'M'.} \\
\\ 
\texttt{Łącznie zgromadziliście 1 z 16 kar.} \\
\texttt{Spróbuj ponownie!} \\
            }
        }
    }
    & & \\ \\

    \theutterance \stepcounter{utterance}  
    & & \multicolumn{4}{p{0.6\linewidth}}{
        \cellcolor[rgb]{0.9,0.9,0.9}{
            \makecell[{{p{\linewidth}}}]{
                \texttt{\tiny{[GM$\rangle$P2]}}
                \texttt{Kara: (5,7) nie jest puste, ale zawiera 'M'.} \\
\\ 
\texttt{Łącznie zgromadziliście 1 z 16 kar.} \\
\texttt{Spróbuj ponownie!} \\
            }
        }
    }
    & \\ \\

    \theutterance \stepcounter{utterance}  
    & & & \multicolumn{4}{p{0.6\linewidth}}{
        \cellcolor[rgb]{0.9,0.9,0.9}{
            \makecell[{{p{\linewidth}}}]{
                \texttt{\tiny{[GM$\langle$P2]}}
                \texttt{RUCH: O, (6,7)} \\
            }
        }
    }
     \\ \\

    \theutterance \stepcounter{utterance}  
    & & & \multicolumn{2}{p{0.3\linewidth}}{
        \cellcolor[rgb]{0.9,0.9,0.9}{
            \makecell[{{p{\linewidth}}}]{
                \texttt{\tiny{[GM$|$GM]}}
                \texttt{RUCH: O, (6,7)} \\
            }
        }
    }
    & & \\ \\

    \theutterance \stepcounter{utterance}  
    & & & \multicolumn{2}{p{0.3\linewidth}}{
        \cellcolor[rgb]{0.9,0.9,0.9}{
            \makecell[{{p{\linewidth}}}]{
                \texttt{\tiny{[GM$|$GM]}}
                \texttt{Kara: (6,7) nie jest puste, ale zawiera '│'.} \\
\\ 
\texttt{Łącznie zgromadziliście 2 z 16 kar.} \\
\texttt{Spróbuj ponownie!} \\
            }
        }
    }
    & & \\ \\

    \theutterance \stepcounter{utterance}  
    & & \multicolumn{4}{p{0.6\linewidth}}{
        \cellcolor[rgb]{0.9,0.9,0.9}{
            \makecell[{{p{\linewidth}}}]{
                \texttt{\tiny{[GM$\rangle$P2]}}
                \texttt{Kara: (6,7) nie jest puste, ale zawiera '│'.} \\
\\ 
\texttt{Łącznie zgromadziliście 2 z 16 kar.} \\
\texttt{Spróbuj ponownie!} \\
            }
        }
    }
    & \\ \\

    \theutterance \stepcounter{utterance}  
    & & & \multicolumn{4}{p{0.6\linewidth}}{
        \cellcolor[rgb]{0.9,0.9,0.9}{
            \makecell[{{p{\linewidth}}}]{
                \texttt{\tiny{[GM$\langle$P2]}}
                \texttt{RUCH: O, (1,7)} \\
            }
        }
    }
     \\ \\

    \theutterance \stepcounter{utterance}  
    & & & \multicolumn{2}{p{0.3\linewidth}}{
        \cellcolor[rgb]{0.9,0.9,0.9}{
            \makecell[{{p{\linewidth}}}]{
                \texttt{\tiny{[GM$|$GM]}}
                \texttt{RUCH: O, (1,7)} \\
            }
        }
    }
    & & \\ \\

    \theutterance \stepcounter{utterance}  
    & & & \multicolumn{2}{p{0.3\linewidth}}{
        \cellcolor[rgb]{0.9,0.9,0.9}{
            \makecell[{{p{\linewidth}}}]{
                \texttt{\tiny{[GM$|$GM]}}
                \texttt{O moved to (1, 7)} \\
            }
        }
    }
    & & \\ \\

    \theutterance \stepcounter{utterance}  
    & & \multicolumn{4}{p{0.6\linewidth}}{
        \cellcolor[rgb]{0.9,0.9,0.9}{
            \makecell[{{p{\linewidth}}}]{
                \texttt{\tiny{[P1$\langle$GM]}}
                \texttt{Twoja wiadomość została przekazana drugiemu graczowi.} \\
\\ 
\\ 
\texttt{Obecnie grasz rundę 2 z maksymalnie 28.} \\
\texttt{Łącznie zgromadziliście 2 z 16 kar.} \\
\texttt{Drugi gracz przesunął obiekt na swojej siatce.} \\
\\ 
\\ 
\texttt{Jakie jest twoje następne polecenie?} \\
            }
        }
    }
    & \\ \\

    \theutterance \stepcounter{utterance}  
    & \multicolumn{4}{p{0.6\linewidth}}{
        \cellcolor[rgb]{0.9,0.9,0.9}{
            \makecell[{{p{\linewidth}}}]{
                \texttt{\tiny{[P1$\rangle$GM]}}
                \texttt{RUCH: P, (1, 7)} \\
            }
        }
    }
    & & \\ \\

    \theutterance \stepcounter{utterance}  
    & & & \multicolumn{2}{p{0.3\linewidth}}{
        \cellcolor[rgb]{0.9,0.9,0.9}{
            \makecell[{{p{\linewidth}}}]{
                \texttt{\tiny{[GM$|$GM]}}
                \texttt{RUCH: P, (1, 7)} \\
            }
        }
    }
    & & \\ \\

    \theutterance \stepcounter{utterance}  
    & & & \multicolumn{2}{p{0.3\linewidth}}{
        \cellcolor[rgb]{0.9,0.9,0.9}{
            \makecell[{{p{\linewidth}}}]{
                \texttt{\tiny{[GM$|$GM]}}
                \texttt{P moved to (1, 7)} \\
            }
        }
    }
    & & \\ \\

    \theutterance \stepcounter{utterance}  
    & & \multicolumn{4}{p{0.6\linewidth}}{
        \cellcolor[rgb]{0.9,0.9,0.9}{
            \makecell[{{p{\linewidth}}}]{
                \texttt{\tiny{[GM$\rangle$P2]}}
                \texttt{Pomyślnie przeniesiono 'O' do (1,7). Twoja zaktualizowana siatka wygląda następująco:} \\
\texttt{```} \\
\texttt{ 1234567} \\
\texttt{╔══╤════╗} \\
\texttt{║T◌│D◌◌A║ 1} \\
\texttt{╟──┤◌◌┌─╢ 2} \\
\texttt{║◌◌├──┤◌║ 3} \\
\texttt{╟──┘◌◌├─╢ 4} \\
\texttt{║◌P◌◌◌│S║ 5} \\
\texttt{╟─┐◌┌─┼─╢ 6} \\
\texttt{║O│◌│M│◌║ 7} \\
\texttt{╚═╧═╧═╧═╝} \\
\texttt{```} \\
\texttt{Obecnie grasz rundę 3 z maksymalnie 28.} \\
\texttt{Łącznie zgromadziliście 2 z 16 kar.} \\
\texttt{Drugi gracz przesunął obiekt na swojej siatce.} \\
\\ 
\\ 
\texttt{Jakie jest twoje następne polecenie?} \\
            }
        }
    }
    & \\ \\

    \theutterance \stepcounter{utterance}  
    & & & \multicolumn{4}{p{0.6\linewidth}}{
        \cellcolor[rgb]{0.9,0.9,0.9}{
            \makecell[{{p{\linewidth}}}]{
                \texttt{\tiny{[GM$\langle$P2]}}
                \texttt{RUCH: D, (7,2)} \\
            }
        }
    }
     \\ \\

    \theutterance \stepcounter{utterance}  
    & & & \multicolumn{2}{p{0.3\linewidth}}{
        \cellcolor[rgb]{0.9,0.9,0.9}{
            \makecell[{{p{\linewidth}}}]{
                \texttt{\tiny{[GM$|$GM]}}
                \texttt{RUCH: D, (7,2)} \\
            }
        }
    }
    & & \\ \\

    \theutterance \stepcounter{utterance}  
    & & & \multicolumn{2}{p{0.3\linewidth}}{
        \cellcolor[rgb]{0.9,0.9,0.9}{
            \makecell[{{p{\linewidth}}}]{
                \texttt{\tiny{[GM$|$GM]}}
                \texttt{Kara: (7,2) nie jest puste, ale zawiera '─'.} \\
\\ 
\texttt{Łącznie zgromadziliście 3 z 16 kar.} \\
\texttt{Spróbuj ponownie!} \\
            }
        }
    }
    & & \\ \\

    \theutterance \stepcounter{utterance}  
    & & \multicolumn{4}{p{0.6\linewidth}}{
        \cellcolor[rgb]{0.9,0.9,0.9}{
            \makecell[{{p{\linewidth}}}]{
                \texttt{\tiny{[GM$\rangle$P2]}}
                \texttt{Kara: (7,2) nie jest puste, ale zawiera '─'.} \\
\\ 
\texttt{Łącznie zgromadziliście 3 z 16 kar.} \\
\texttt{Spróbuj ponownie!} \\
            }
        }
    }
    & \\ \\

    \theutterance \stepcounter{utterance}  
    & & & \multicolumn{4}{p{0.6\linewidth}}{
        \cellcolor[rgb]{0.9,0.9,0.9}{
            \makecell[{{p{\linewidth}}}]{
                \texttt{\tiny{[GM$\langle$P2]}}
                \texttt{RUCH: D, (7,3)} \\
            }
        }
    }
     \\ \\

    \theutterance \stepcounter{utterance}  
    & & & \multicolumn{2}{p{0.3\linewidth}}{
        \cellcolor[rgb]{0.9,0.9,0.9}{
            \makecell[{{p{\linewidth}}}]{
                \texttt{\tiny{[GM$|$GM]}}
                \texttt{RUCH: D, (7,3)} \\
            }
        }
    }
    & & \\ \\

    \theutterance \stepcounter{utterance}  
    & & & \multicolumn{2}{p{0.3\linewidth}}{
        \cellcolor[rgb]{0.9,0.9,0.9}{
            \makecell[{{p{\linewidth}}}]{
                \texttt{\tiny{[GM$|$GM]}}
                \texttt{D moved to (7, 3)} \\
            }
        }
    }
    & & \\ \\

    \theutterance \stepcounter{utterance}  
    & & \multicolumn{4}{p{0.6\linewidth}}{
        \cellcolor[rgb]{0.9,0.9,0.9}{
            \makecell[{{p{\linewidth}}}]{
                \texttt{\tiny{[P1$\langle$GM]}}
                \texttt{Pomyślnie przeniesiono 'P' do (1,7). Twoja zaktualizowana siatka wygląda następująco:} \\
\texttt{```} \\
\texttt{ 1234567} \\
\texttt{╔══╤════╗} \\
\texttt{║◌◌│◌◌◌A║ 1} \\
\texttt{╟──┤◌◌┌─╢ 2} \\
\texttt{║◌◌├──┤T║ 3} \\
\texttt{╟──┘M◌├─╢ 4} \\
\texttt{║◌◌OD◌│S║ 5} \\
\texttt{╟─┐◌┌─┼─╢ 6} \\
\texttt{║P│◌│◌│◌║ 7} \\
\texttt{╚═╧═╧═╧═╝} \\
\texttt{```} \\
\texttt{Obecnie grasz rundę 3 z maksymalnie 28.} \\
\texttt{Łącznie zgromadziliście 3 z 16 kar.} \\
\texttt{Drugi gracz przesunął obiekt na swojej siatce.} \\
\\ 
\\ 
\texttt{Jakie jest twoje następne polecenie?} \\
            }
        }
    }
    & \\ \\

    \theutterance \stepcounter{utterance}  
    & \multicolumn{4}{p{0.6\linewidth}}{
        \cellcolor[rgb]{0.9,0.9,0.9}{
            \makecell[{{p{\linewidth}}}]{
                \texttt{\tiny{[P1$\rangle$GM]}}
                \texttt{RUCH: D, (7, 2)} \\
            }
        }
    }
    & & \\ \\

    \theutterance \stepcounter{utterance}  
    & & & \multicolumn{2}{p{0.3\linewidth}}{
        \cellcolor[rgb]{0.9,0.9,0.9}{
            \makecell[{{p{\linewidth}}}]{
                \texttt{\tiny{[GM$|$GM]}}
                \texttt{RUCH: D, (7, 2)} \\
            }
        }
    }
    & & \\ \\

    \theutterance \stepcounter{utterance}  
    & & & \multicolumn{2}{p{0.3\linewidth}}{
        \cellcolor[rgb]{0.9,0.9,0.9}{
            \makecell[{{p{\linewidth}}}]{
                \texttt{\tiny{[GM$|$GM]}}
                \texttt{Kara: (7,2) nie jest puste, ale zawiera '─'.} \\
\\ 
\texttt{Łącznie zgromadziliście 4 z 16 kar.} \\
\texttt{Spróbuj ponownie!} \\
            }
        }
    }
    & & \\ \\

    \theutterance \stepcounter{utterance}  
    & & \multicolumn{4}{p{0.6\linewidth}}{
        \cellcolor[rgb]{0.9,0.9,0.9}{
            \makecell[{{p{\linewidth}}}]{
                \texttt{\tiny{[P1$\langle$GM]}}
                \texttt{Kara: (7,2) nie jest puste, ale zawiera '─'.} \\
\\ 
\texttt{Łącznie zgromadziliście 4 z 16 kar.} \\
\texttt{Spróbuj ponownie!} \\
            }
        }
    }
    & \\ \\

    \theutterance \stepcounter{utterance}  
    & \multicolumn{4}{p{0.6\linewidth}}{
        \cellcolor[rgb]{0.9,0.9,0.9}{
            \makecell[{{p{\linewidth}}}]{
                \texttt{\tiny{[P1$\rangle$GM]}}
                \texttt{RUCH: D, (6, 2)} \\
            }
        }
    }
    & & \\ \\

    \theutterance \stepcounter{utterance}  
    & & & \multicolumn{2}{p{0.3\linewidth}}{
        \cellcolor[rgb]{0.9,0.9,0.9}{
            \makecell[{{p{\linewidth}}}]{
                \texttt{\tiny{[GM$|$GM]}}
                \texttt{RUCH: D, (6, 2)} \\
            }
        }
    }
    & & \\ \\

    \theutterance \stepcounter{utterance}  
    & & & \multicolumn{2}{p{0.3\linewidth}}{
        \cellcolor[rgb]{0.9,0.9,0.9}{
            \makecell[{{p{\linewidth}}}]{
                \texttt{\tiny{[GM$|$GM]}}
                \texttt{Kara: (6,2) nie jest puste, ale zawiera '┌'.} \\
\\ 
\texttt{Łącznie zgromadziliście 5 z 16 kar.} \\
\texttt{Spróbuj ponownie!} \\
            }
        }
    }
    & & \\ \\

    \theutterance \stepcounter{utterance}  
    & & \multicolumn{4}{p{0.6\linewidth}}{
        \cellcolor[rgb]{0.9,0.9,0.9}{
            \makecell[{{p{\linewidth}}}]{
                \texttt{\tiny{[P1$\langle$GM]}}
                \texttt{Kara: (6,2) nie jest puste, ale zawiera '┌'.} \\
\\ 
\texttt{Łącznie zgromadziliście 5 z 16 kar.} \\
\texttt{Spróbuj ponownie!} \\
            }
        }
    }
    & \\ \\

    \theutterance \stepcounter{utterance}  
    & \multicolumn{4}{p{0.6\linewidth}}{
        \cellcolor[rgb]{0.9,0.9,0.9}{
            \makecell[{{p{\linewidth}}}]{
                \texttt{\tiny{[P1$\rangle$GM]}}
                \texttt{RUCH: D, (5, 2)} \\
            }
        }
    }
    & & \\ \\

    \theutterance \stepcounter{utterance}  
    & & & \multicolumn{2}{p{0.3\linewidth}}{
        \cellcolor[rgb]{0.9,0.9,0.9}{
            \makecell[{{p{\linewidth}}}]{
                \texttt{\tiny{[GM$|$GM]}}
                \texttt{RUCH: D, (5, 2)} \\
            }
        }
    }
    & & \\ \\

    \theutterance \stepcounter{utterance}  
    & & & \multicolumn{2}{p{0.3\linewidth}}{
        \cellcolor[rgb]{0.9,0.9,0.9}{
            \makecell[{{p{\linewidth}}}]{
                \texttt{\tiny{[GM$|$GM]}}
                \texttt{D moved to (5, 2)} \\
            }
        }
    }
    & & \\ \\

    \theutterance \stepcounter{utterance}  
    & & \multicolumn{4}{p{0.6\linewidth}}{
        \cellcolor[rgb]{0.9,0.9,0.9}{
            \makecell[{{p{\linewidth}}}]{
                \texttt{\tiny{[GM$\rangle$P2]}}
                \texttt{Pomyślnie przeniesiono 'D' do (7,3). Twoja zaktualizowana siatka wygląda następująco:} \\
\texttt{```} \\
\texttt{ 1234567} \\
\texttt{╔══╤════╗} \\
\texttt{║T◌│◌◌◌A║ 1} \\
\texttt{╟──┤◌◌┌─╢ 2} \\
\texttt{║◌◌├──┤D║ 3} \\
\texttt{╟──┘◌◌├─╢ 4} \\
\texttt{║◌P◌◌◌│S║ 5} \\
\texttt{╟─┐◌┌─┼─╢ 6} \\
\texttt{║O│◌│M│◌║ 7} \\
\texttt{╚═╧═╧═╧═╝} \\
\texttt{```} \\
\texttt{Obecnie grasz rundę 4 z maksymalnie 28.} \\
\texttt{Łącznie zgromadziliście 5 z 16 kar.} \\
\texttt{Drugi gracz przesunął obiekt na swojej siatce.} \\
\\ 
\\ 
\texttt{Jakie jest twoje następne polecenie?} \\
            }
        }
    }
    & \\ \\

    \theutterance \stepcounter{utterance}  
    & & & \multicolumn{4}{p{0.6\linewidth}}{
        \cellcolor[rgb]{0.9,0.9,0.9}{
            \makecell[{{p{\linewidth}}}]{
                \texttt{\tiny{[GM$\langle$P2]}}
                \texttt{POWIEDZ: D poprawiłem na (7,3), bo (7,2) zawiera krawędź. Reszta bez zmian: A=(6,1), T=(3,1), M=(5,1), S=(7,1), O=(1,7), P=(6,7).} \\
            }
        }
    }
     \\ \\

    \theutterance \stepcounter{utterance}  
    & & & \multicolumn{2}{p{0.3\linewidth}}{
        \cellcolor[rgb]{0.9,0.9,0.9}{
            \makecell[{{p{\linewidth}}}]{
                \texttt{\tiny{[GM$|$GM]}}
                \texttt{POWIEDZ: D poprawiłem na (7,3), bo (7,2) zawiera krawędź. Reszta bez zmian: A=(6,1), T=(3,1), M=(5,1), S=(7,1), O=(1,7), P=(6,7).} \\
            }
        }
    }
    & & \\ \\

    \theutterance \stepcounter{utterance}  
    & & \multicolumn{4}{p{0.6\linewidth}}{
        \cellcolor[rgb]{0.9,0.9,0.9}{
            \makecell[{{p{\linewidth}}}]{
                \texttt{\tiny{[P1$\langle$GM]}}
                \texttt{Pomyślnie przeniesiono 'D' do (5,2). Twoja zaktualizowana siatka wygląda następująco:} \\
\texttt{```} \\
\texttt{ 1234567} \\
\texttt{╔══╤════╗} \\
\texttt{║◌◌│◌◌◌A║ 1} \\
\texttt{╟──┤◌D┌─╢ 2} \\
\texttt{║◌◌├──┤T║ 3} \\
\texttt{╟──┘M◌├─╢ 4} \\
\texttt{║◌◌O◌◌│S║ 5} \\
\texttt{╟─┐◌┌─┼─╢ 6} \\
\texttt{║P│◌│◌│◌║ 7} \\
\texttt{╚═╧═╧═╧═╝} \\
\texttt{```} \\
\texttt{Obecnie grasz rundę 4 z maksymalnie 28.} \\
\texttt{Łącznie zgromadziliście 5 z 16 kar.} \\
\texttt{Drugi gracz wysłał tę wiadomość:} \\
\\ 
\\ 
\texttt{"D poprawiłem na (7,3), bo (7,2) zawiera krawędź. Reszta bez zmian: A=(6,1), T=(3,1), M=(5,1), S=(7,1), O=(1,7), P=(6,7)."} \\
\\ 
\\ 
\texttt{Jakie jest twoje następne polecenie?} \\
            }
        }
    }
    & \\ \\

    \theutterance \stepcounter{utterance}  
    & \multicolumn{4}{p{0.6\linewidth}}{
        \cellcolor[rgb]{0.9,0.9,0.9}{
            \makecell[{{p{\linewidth}}}]{
                \texttt{\tiny{[P1$\rangle$GM]}}
                \texttt{POWIEDZ: D zrozumiałem, zostawię go u siebie na (5,2) w ramach kompromisu. Mój nowy cel to: A=(6,1), T=(3,1), M=(5,1), S=(7,1), O=(1,7), P=(6,7), D=(5,2). Zaczynam przesuwać O z (3,5) do (1,7).} \\
            }
        }
    }
    & & \\ \\

    \theutterance \stepcounter{utterance}  
    & & & \multicolumn{2}{p{0.3\linewidth}}{
        \cellcolor[rgb]{0.9,0.9,0.9}{
            \makecell[{{p{\linewidth}}}]{
                \texttt{\tiny{[GM$|$GM]}}
                \texttt{POWIEDZ: D zrozumiałem, zostawię go u siebie na (5,2) w ramach kompromisu. Mój nowy cel to: A=(6,1), T=(3,1), M=(5,1), S=(7,1), O=(1,7), P=(6,7), D=(5,2). Zaczynam przesuwać O z (3,5) do (1,7).} \\
            }
        }
    }
    & & \\ \\

    \theutterance \stepcounter{utterance}  
    & & \multicolumn{4}{p{0.6\linewidth}}{
        \cellcolor[rgb]{0.9,0.9,0.9}{
            \makecell[{{p{\linewidth}}}]{
                \texttt{\tiny{[GM$\rangle$P2]}}
                \texttt{Twoja wiadomość została przekazana drugiemu graczowi.} \\
\\ 
\\ 
\texttt{Obecnie grasz rundę 5 z maksymalnie 28.} \\
\texttt{Łącznie zgromadziliście 5 z 16 kar.} \\
\texttt{Drugi gracz wysłał tę wiadomość:} \\
\\ 
\\ 
\texttt{"D zrozumiałem, zostawię go u siebie na (5,2) w ramach kompromisu. Mój nowy cel to: A=(6,1), T=(3,1), M=(5,1), S=(7,1), O=(1,7), P=(6,7), D=(5,2). Zaczynam przesuwać O z (3,5) do (1,7)."} \\
\\ 
\\ 
\texttt{Jakie jest twoje następne polecenie?} \\
            }
        }
    }
    & \\ \\

    \theutterance \stepcounter{utterance}  
    & & & \multicolumn{4}{p{0.6\linewidth}}{
        \cellcolor[rgb]{0.9,0.9,0.9}{
            \makecell[{{p{\linewidth}}}]{
                \texttt{\tiny{[GM$\langle$P2]}}
                \texttt{RUCH: T, (3,1)} \\
            }
        }
    }
     \\ \\

    \theutterance \stepcounter{utterance}  
    & & & \multicolumn{2}{p{0.3\linewidth}}{
        \cellcolor[rgb]{0.9,0.9,0.9}{
            \makecell[{{p{\linewidth}}}]{
                \texttt{\tiny{[GM$|$GM]}}
                \texttt{RUCH: T, (3,1)} \\
            }
        }
    }
    & & \\ \\

    \theutterance \stepcounter{utterance}  
    & & & \multicolumn{2}{p{0.3\linewidth}}{
        \cellcolor[rgb]{0.9,0.9,0.9}{
            \makecell[{{p{\linewidth}}}]{
                \texttt{\tiny{[GM$|$GM]}}
                \texttt{Kara: (3,1) nie jest puste, ale zawiera '│'.} \\
\\ 
\texttt{Łącznie zgromadziliście 6 z 16 kar.} \\
\texttt{Spróbuj ponownie!} \\
            }
        }
    }
    & & \\ \\

    \theutterance \stepcounter{utterance}  
    & & \multicolumn{4}{p{0.6\linewidth}}{
        \cellcolor[rgb]{0.9,0.9,0.9}{
            \makecell[{{p{\linewidth}}}]{
                \texttt{\tiny{[GM$\rangle$P2]}}
                \texttt{Kara: (3,1) nie jest puste, ale zawiera '│'.} \\
\\ 
\texttt{Łącznie zgromadziliście 6 z 16 kar.} \\
\texttt{Spróbuj ponownie!} \\
            }
        }
    }
    & \\ \\

    \theutterance \stepcounter{utterance}  
    & & & \multicolumn{4}{p{0.6\linewidth}}{
        \cellcolor[rgb]{0.9,0.9,0.9}{
            \makecell[{{p{\linewidth}}}]{
                \texttt{\tiny{[GM$\langle$P2]}}
                \texttt{RUCH: T, (2,1)} \\
            }
        }
    }
     \\ \\

    \theutterance \stepcounter{utterance}  
    & & & \multicolumn{2}{p{0.3\linewidth}}{
        \cellcolor[rgb]{0.9,0.9,0.9}{
            \makecell[{{p{\linewidth}}}]{
                \texttt{\tiny{[GM$|$GM]}}
                \texttt{RUCH: T, (2,1)} \\
            }
        }
    }
    & & \\ \\

    \theutterance \stepcounter{utterance}  
    & & & \multicolumn{2}{p{0.3\linewidth}}{
        \cellcolor[rgb]{0.9,0.9,0.9}{
            \makecell[{{p{\linewidth}}}]{
                \texttt{\tiny{[GM$|$GM]}}
                \texttt{T moved to (2, 1)} \\
            }
        }
    }
    & & \\ \\

    \theutterance \stepcounter{utterance}  
    & & \multicolumn{4}{p{0.6\linewidth}}{
        \cellcolor[rgb]{0.9,0.9,0.9}{
            \makecell[{{p{\linewidth}}}]{
                \texttt{\tiny{[P1$\langle$GM]}}
                \texttt{Twoja wiadomość została przekazana drugiemu graczowi.} \\
\\ 
\\ 
\texttt{Obecnie grasz rundę 5 z maksymalnie 28.} \\
\texttt{Łącznie zgromadziliście 6 z 16 kar.} \\
\texttt{Drugi gracz przesunął obiekt na swojej siatce.} \\
\\ 
\\ 
\texttt{Jakie jest twoje następne polecenie?} \\
            }
        }
    }
    & \\ \\

    \theutterance \stepcounter{utterance}  
    & \multicolumn{4}{p{0.6\linewidth}}{
        \cellcolor[rgb]{0.9,0.9,0.9}{
            \makecell[{{p{\linewidth}}}]{
                \texttt{\tiny{[P1$\rangle$GM]}}
                \texttt{RUCH: O, (2, 5)} \\
            }
        }
    }
    & & \\ \\

    \theutterance \stepcounter{utterance}  
    & & & \multicolumn{2}{p{0.3\linewidth}}{
        \cellcolor[rgb]{0.9,0.9,0.9}{
            \makecell[{{p{\linewidth}}}]{
                \texttt{\tiny{[GM$|$GM]}}
                \texttt{RUCH: O, (2, 5)} \\
            }
        }
    }
    & & \\ \\

    \theutterance \stepcounter{utterance}  
    & & & \multicolumn{2}{p{0.3\linewidth}}{
        \cellcolor[rgb]{0.9,0.9,0.9}{
            \makecell[{{p{\linewidth}}}]{
                \texttt{\tiny{[GM$|$GM]}}
                \texttt{O moved to (2, 5)} \\
            }
        }
    }
    & & \\ \\

    \theutterance \stepcounter{utterance}  
    & & \multicolumn{4}{p{0.6\linewidth}}{
        \cellcolor[rgb]{0.9,0.9,0.9}{
            \makecell[{{p{\linewidth}}}]{
                \texttt{\tiny{[GM$\rangle$P2]}}
                \texttt{Pomyślnie przeniesiono 'T' do (2,1). Twoja zaktualizowana siatka wygląda następująco:} \\
\texttt{```} \\
\texttt{ 1234567} \\
\texttt{╔══╤════╗} \\
\texttt{║◌T│◌◌◌A║ 1} \\
\texttt{╟──┤◌◌┌─╢ 2} \\
\texttt{║◌◌├──┤D║ 3} \\
\texttt{╟──┘◌◌├─╢ 4} \\
\texttt{║◌P◌◌◌│S║ 5} \\
\texttt{╟─┐◌┌─┼─╢ 6} \\
\texttt{║O│◌│M│◌║ 7} \\
\texttt{╚═╧═╧═╧═╝} \\
\texttt{```} \\
\texttt{Obecnie grasz rundę 6 z maksymalnie 28.} \\
\texttt{Łącznie zgromadziliście 6 z 16 kar.} \\
\texttt{Drugi gracz przesunął obiekt na swojej siatce.} \\
\\ 
\\ 
\texttt{Jakie jest twoje następne polecenie?} \\
            }
        }
    }
    & \\ \\

    \theutterance \stepcounter{utterance}  
    & & & \multicolumn{4}{p{0.6\linewidth}}{
        \cellcolor[rgb]{0.9,0.9,0.9}{
            \makecell[{{p{\linewidth}}}]{
                \texttt{\tiny{[GM$\langle$P2]}}
                \texttt{POWIEDZ: T poprawiłem na (2,1), bo (3,1) zawiera krawędź. Docelowo: A=(6,1), T=(2,1), M=(5,1), S=(7,1), O=(1,7), P=(6,7), D=(5,2).} \\
            }
        }
    }
     \\ \\

    \theutterance \stepcounter{utterance}  
    & & & \multicolumn{2}{p{0.3\linewidth}}{
        \cellcolor[rgb]{0.9,0.9,0.9}{
            \makecell[{{p{\linewidth}}}]{
                \texttt{\tiny{[GM$|$GM]}}
                \texttt{POWIEDZ: T poprawiłem na (2,1), bo (3,1) zawiera krawędź. Docelowo: A=(6,1), T=(2,1), M=(5,1), S=(7,1), O=(1,7), P=(6,7), D=(5,2).} \\
            }
        }
    }
    & & \\ \\

    \theutterance \stepcounter{utterance}  
    & & \multicolumn{4}{p{0.6\linewidth}}{
        \cellcolor[rgb]{0.9,0.9,0.9}{
            \makecell[{{p{\linewidth}}}]{
                \texttt{\tiny{[P1$\langle$GM]}}
                \texttt{Pomyślnie przeniesiono 'O' do (2,5). Twoja zaktualizowana siatka wygląda następująco:} \\
\texttt{```} \\
\texttt{ 1234567} \\
\texttt{╔══╤════╗} \\
\texttt{║◌◌│◌◌◌A║ 1} \\
\texttt{╟──┤◌D┌─╢ 2} \\
\texttt{║◌◌├──┤T║ 3} \\
\texttt{╟──┘M◌├─╢ 4} \\
\texttt{║◌O◌◌◌│S║ 5} \\
\texttt{╟─┐◌┌─┼─╢ 6} \\
\texttt{║P│◌│◌│◌║ 7} \\
\texttt{╚═╧═╧═╧═╝} \\
\texttt{```} \\
\texttt{Obecnie grasz rundę 6 z maksymalnie 28.} \\
\texttt{Łącznie zgromadziliście 6 z 16 kar.} \\
\texttt{Drugi gracz wysłał tę wiadomość:} \\
\\ 
\\ 
\texttt{"T poprawiłem na (2,1), bo (3,1) zawiera krawędź. Docelowo: A=(6,1), T=(2,1), M=(5,1), S=(7,1), O=(1,7), P=(6,7), D=(5,2)."} \\
\\ 
\\ 
\texttt{Jakie jest twoje następne polecenie?} \\
            }
        }
    }
    & \\ \\

    \theutterance \stepcounter{utterance}  
    & \multicolumn{4}{p{0.6\linewidth}}{
        \cellcolor[rgb]{0.9,0.9,0.9}{
            \makecell[{{p{\linewidth}}}]{
                \texttt{\tiny{[P1$\rangle$GM]}}
                \texttt{POWIEDZ: Zrozumiałem, T=(2,1) jest u mnie puste i akceptuję ten cel. Przesuwam O z (2,5) do (1,5).} \\
            }
        }
    }
    & & \\ \\

    \theutterance \stepcounter{utterance}  
    & & & \multicolumn{2}{p{0.3\linewidth}}{
        \cellcolor[rgb]{0.9,0.9,0.9}{
            \makecell[{{p{\linewidth}}}]{
                \texttt{\tiny{[GM$|$GM]}}
                \texttt{POWIEDZ: Zrozumiałem, T=(2,1) jest u mnie puste i akceptuję ten cel. Przesuwam O z (2,5) do (1,5).} \\
            }
        }
    }
    & & \\ \\

    \theutterance \stepcounter{utterance}  
    & & \multicolumn{4}{p{0.6\linewidth}}{
        \cellcolor[rgb]{0.9,0.9,0.9}{
            \makecell[{{p{\linewidth}}}]{
                \texttt{\tiny{[GM$\rangle$P2]}}
                \texttt{Twoja wiadomość została przekazana drugiemu graczowi.} \\
\\ 
\\ 
\texttt{Obecnie grasz rundę 7 z maksymalnie 28.} \\
\texttt{Łącznie zgromadziliście 6 z 16 kar.} \\
\texttt{Drugi gracz wysłał tę wiadomość:} \\
\\ 
\\ 
\texttt{"Zrozumiałem, T=(2,1) jest u mnie puste i akceptuję ten cel. Przesuwam O z (2,5) do (1,5)."} \\
\\ 
\\ 
\texttt{Jakie jest twoje następne polecenie?} \\
            }
        }
    }
    & \\ \\

    \theutterance \stepcounter{utterance}  
    & & & \multicolumn{4}{p{0.6\linewidth}}{
        \cellcolor[rgb]{0.9,0.9,0.9}{
            \makecell[{{p{\linewidth}}}]{
                \texttt{\tiny{[GM$\langle$P2]}}
                \texttt{RUCH: M, (5,1)} \\
            }
        }
    }
     \\ \\

    \theutterance \stepcounter{utterance}  
    & & & \multicolumn{2}{p{0.3\linewidth}}{
        \cellcolor[rgb]{0.9,0.9,0.9}{
            \makecell[{{p{\linewidth}}}]{
                \texttt{\tiny{[GM$|$GM]}}
                \texttt{RUCH: M, (5,1)} \\
            }
        }
    }
    & & \\ \\

    \theutterance \stepcounter{utterance}  
    & & & \multicolumn{2}{p{0.3\linewidth}}{
        \cellcolor[rgb]{0.9,0.9,0.9}{
            \makecell[{{p{\linewidth}}}]{
                \texttt{\tiny{[GM$|$GM]}}
                \texttt{M moved to (5, 1)} \\
            }
        }
    }
    & & \\ \\

    \theutterance \stepcounter{utterance}  
    & & \multicolumn{4}{p{0.6\linewidth}}{
        \cellcolor[rgb]{0.9,0.9,0.9}{
            \makecell[{{p{\linewidth}}}]{
                \texttt{\tiny{[P1$\langle$GM]}}
                \texttt{Twoja wiadomość została przekazana drugiemu graczowi.} \\
\\ 
\\ 
\texttt{Obecnie grasz rundę 7 z maksymalnie 28.} \\
\texttt{Łącznie zgromadziliście 6 z 16 kar.} \\
\texttt{Drugi gracz przesunął obiekt na swojej siatce.} \\
\\ 
\\ 
\texttt{Jakie jest twoje następne polecenie?} \\
            }
        }
    }
    & \\ \\

    \theutterance \stepcounter{utterance}  
    & \multicolumn{4}{p{0.6\linewidth}}{
        \cellcolor[rgb]{0.9,0.9,0.9}{
            \makecell[{{p{\linewidth}}}]{
                \texttt{\tiny{[P1$\rangle$GM]}}
                \texttt{RUCH: O, (1, 5)} \\
            }
        }
    }
    & & \\ \\

    \theutterance \stepcounter{utterance}  
    & & & \multicolumn{2}{p{0.3\linewidth}}{
        \cellcolor[rgb]{0.9,0.9,0.9}{
            \makecell[{{p{\linewidth}}}]{
                \texttt{\tiny{[GM$|$GM]}}
                \texttt{RUCH: O, (1, 5)} \\
            }
        }
    }
    & & \\ \\

    \theutterance \stepcounter{utterance}  
    & & & \multicolumn{2}{p{0.3\linewidth}}{
        \cellcolor[rgb]{0.9,0.9,0.9}{
            \makecell[{{p{\linewidth}}}]{
                \texttt{\tiny{[GM$|$GM]}}
                \texttt{O moved to (1, 5)} \\
            }
        }
    }
    & & \\ \\

    \theutterance \stepcounter{utterance}  
    & & \multicolumn{4}{p{0.6\linewidth}}{
        \cellcolor[rgb]{0.9,0.9,0.9}{
            \makecell[{{p{\linewidth}}}]{
                \texttt{\tiny{[GM$\rangle$P2]}}
                \texttt{Pomyślnie przeniesiono 'M' do (5,1). Twoja zaktualizowana siatka wygląda następująco:} \\
\texttt{```} \\
\texttt{ 1234567} \\
\texttt{╔══╤════╗} \\
\texttt{║◌T│◌M◌A║ 1} \\
\texttt{╟──┤◌◌┌─╢ 2} \\
\texttt{║◌◌├──┤D║ 3} \\
\texttt{╟──┘◌◌├─╢ 4} \\
\texttt{║◌P◌◌◌│S║ 5} \\
\texttt{╟─┐◌┌─┼─╢ 6} \\
\texttt{║O│◌│◌│◌║ 7} \\
\texttt{╚═╧═╧═╧═╝} \\
\texttt{```} \\
\texttt{Obecnie grasz rundę 8 z maksymalnie 28.} \\
\texttt{Łącznie zgromadziliście 6 z 16 kar.} \\
\texttt{Drugi gracz przesunął obiekt na swojej siatce.} \\
\\ 
\\ 
\texttt{Jakie jest twoje następne polecenie?} \\
            }
        }
    }
    & \\ \\

    \theutterance \stepcounter{utterance}  
    & & & \multicolumn{4}{p{0.6\linewidth}}{
        \cellcolor[rgb]{0.9,0.9,0.9}{
            \makecell[{{p{\linewidth}}}]{
                \texttt{\tiny{[GM$\langle$P2]}}
                \texttt{RUCH: A, (6,1)} \\
            }
        }
    }
     \\ \\

    \theutterance \stepcounter{utterance}  
    & & & \multicolumn{2}{p{0.3\linewidth}}{
        \cellcolor[rgb]{0.9,0.9,0.9}{
            \makecell[{{p{\linewidth}}}]{
                \texttt{\tiny{[GM$|$GM]}}
                \texttt{RUCH: A, (6,1)} \\
            }
        }
    }
    & & \\ \\

    \theutterance \stepcounter{utterance}  
    & & & \multicolumn{2}{p{0.3\linewidth}}{
        \cellcolor[rgb]{0.9,0.9,0.9}{
            \makecell[{{p{\linewidth}}}]{
                \texttt{\tiny{[GM$|$GM]}}
                \texttt{A moved to (6, 1)} \\
            }
        }
    }
    & & \\ \\

    \theutterance \stepcounter{utterance}  
    & & \multicolumn{4}{p{0.6\linewidth}}{
        \cellcolor[rgb]{0.9,0.9,0.9}{
            \makecell[{{p{\linewidth}}}]{
                \texttt{\tiny{[P1$\langle$GM]}}
                \texttt{Pomyślnie przeniesiono 'O' do (1,5). Twoja zaktualizowana siatka wygląda następująco:} \\
\texttt{```} \\
\texttt{ 1234567} \\
\texttt{╔══╤════╗} \\
\texttt{║◌◌│◌◌◌A║ 1} \\
\texttt{╟──┤◌D┌─╢ 2} \\
\texttt{║◌◌├──┤T║ 3} \\
\texttt{╟──┘M◌├─╢ 4} \\
\texttt{║O◌◌◌◌│S║ 5} \\
\texttt{╟─┐◌┌─┼─╢ 6} \\
\texttt{║P│◌│◌│◌║ 7} \\
\texttt{╚═╧═╧═╧═╝} \\
\texttt{```} \\
\texttt{Obecnie grasz rundę 8 z maksymalnie 28.} \\
\texttt{Łącznie zgromadziliście 6 z 16 kar.} \\
\texttt{Drugi gracz przesunął obiekt na swojej siatce.} \\
\\ 
\\ 
\texttt{Jakie jest twoje następne polecenie?} \\
            }
        }
    }
    & \\ \\

    \theutterance \stepcounter{utterance}  
    & \multicolumn{4}{p{0.6\linewidth}}{
        \cellcolor[rgb]{0.9,0.9,0.9}{
            \makecell[{{p{\linewidth}}}]{
                \texttt{\tiny{[P1$\rangle$GM]}}
                \texttt{RUCH: P, (6, 7)} \\
            }
        }
    }
    & & \\ \\

    \theutterance \stepcounter{utterance}  
    & & & \multicolumn{2}{p{0.3\linewidth}}{
        \cellcolor[rgb]{0.9,0.9,0.9}{
            \makecell[{{p{\linewidth}}}]{
                \texttt{\tiny{[GM$|$GM]}}
                \texttt{RUCH: P, (6, 7)} \\
            }
        }
    }
    & & \\ \\

    \theutterance \stepcounter{utterance}  
    & & & \multicolumn{2}{p{0.3\linewidth}}{
        \cellcolor[rgb]{0.9,0.9,0.9}{
            \makecell[{{p{\linewidth}}}]{
                \texttt{\tiny{[GM$|$GM]}}
                \texttt{Kara: (6,7) nie jest puste, ale zawiera '│'.} \\
\\ 
\texttt{Łącznie zgromadziliście 7 z 16 kar.} \\
\texttt{Spróbuj ponownie!} \\
            }
        }
    }
    & & \\ \\

    \theutterance \stepcounter{utterance}  
    & & \multicolumn{4}{p{0.6\linewidth}}{
        \cellcolor[rgb]{0.9,0.9,0.9}{
            \makecell[{{p{\linewidth}}}]{
                \texttt{\tiny{[P1$\langle$GM]}}
                \texttt{Kara: (6,7) nie jest puste, ale zawiera '│'.} \\
\\ 
\texttt{Łącznie zgromadziliście 7 z 16 kar.} \\
\texttt{Spróbuj ponownie!} \\
            }
        }
    }
    & \\ \\

    \theutterance \stepcounter{utterance}  
    & \multicolumn{4}{p{0.6\linewidth}}{
        \cellcolor[rgb]{0.9,0.9,0.9}{
            \makecell[{{p{\linewidth}}}]{
                \texttt{\tiny{[P1$\rangle$GM]}}
                \texttt{RUCH: P, (5, 7)} \\
            }
        }
    }
    & & \\ \\

    \theutterance \stepcounter{utterance}  
    & & & \multicolumn{2}{p{0.3\linewidth}}{
        \cellcolor[rgb]{0.9,0.9,0.9}{
            \makecell[{{p{\linewidth}}}]{
                \texttt{\tiny{[GM$|$GM]}}
                \texttt{RUCH: P, (5, 7)} \\
            }
        }
    }
    & & \\ \\

    \theutterance \stepcounter{utterance}  
    & & & \multicolumn{2}{p{0.3\linewidth}}{
        \cellcolor[rgb]{0.9,0.9,0.9}{
            \makecell[{{p{\linewidth}}}]{
                \texttt{\tiny{[GM$|$GM]}}
                \texttt{P moved to (5, 7)} \\
            }
        }
    }
    & & \\ \\

    \theutterance \stepcounter{utterance}  
    & & \multicolumn{4}{p{0.6\linewidth}}{
        \cellcolor[rgb]{0.9,0.9,0.9}{
            \makecell[{{p{\linewidth}}}]{
                \texttt{\tiny{[GM$\rangle$P2]}}
                \texttt{Pomyślnie przeniesiono 'A' do (6,1). Twoja zaktualizowana siatka wygląda następująco:} \\
\texttt{```} \\
\texttt{ 1234567} \\
\texttt{╔══╤════╗} \\
\texttt{║◌T│◌MA◌║ 1} \\
\texttt{╟──┤◌◌┌─╢ 2} \\
\texttt{║◌◌├──┤D║ 3} \\
\texttt{╟──┘◌◌├─╢ 4} \\
\texttt{║◌P◌◌◌│S║ 5} \\
\texttt{╟─┐◌┌─┼─╢ 6} \\
\texttt{║O│◌│◌│◌║ 7} \\
\texttt{╚═╧═╧═╧═╝} \\
\texttt{```} \\
\texttt{Obecnie grasz rundę 9 z maksymalnie 28.} \\
\texttt{Łącznie zgromadziliście 7 z 16 kar.} \\
\texttt{Drugi gracz przesunął obiekt na swojej siatce.} \\
\\ 
\\ 
\texttt{Jakie jest twoje następne polecenie?} \\
            }
        }
    }
    & \\ \\

    \theutterance \stepcounter{utterance}  
    & & & \multicolumn{4}{p{0.6\linewidth}}{
        \cellcolor[rgb]{0.9,0.9,0.9}{
            \makecell[{{p{\linewidth}}}]{
                \texttt{\tiny{[GM$\langle$P2]}}
                \texttt{RUCH: S, (7,1)} \\
            }
        }
    }
     \\ \\

    \theutterance \stepcounter{utterance}  
    & & & \multicolumn{2}{p{0.3\linewidth}}{
        \cellcolor[rgb]{0.9,0.9,0.9}{
            \makecell[{{p{\linewidth}}}]{
                \texttt{\tiny{[GM$|$GM]}}
                \texttt{RUCH: S, (7,1)} \\
            }
        }
    }
    & & \\ \\

    \theutterance \stepcounter{utterance}  
    & & & \multicolumn{2}{p{0.3\linewidth}}{
        \cellcolor[rgb]{0.9,0.9,0.9}{
            \makecell[{{p{\linewidth}}}]{
                \texttt{\tiny{[GM$|$GM]}}
                \texttt{S moved to (7, 1)} \\
            }
        }
    }
    & & \\ \\

    \theutterance \stepcounter{utterance}  
    & & \multicolumn{4}{p{0.6\linewidth}}{
        \cellcolor[rgb]{0.9,0.9,0.9}{
            \makecell[{{p{\linewidth}}}]{
                \texttt{\tiny{[P1$\langle$GM]}}
                \texttt{Pomyślnie przeniesiono 'P' do (5,7). Twoja zaktualizowana siatka wygląda następująco:} \\
\texttt{```} \\
\texttt{ 1234567} \\
\texttt{╔══╤════╗} \\
\texttt{║◌◌│◌◌◌A║ 1} \\
\texttt{╟──┤◌D┌─╢ 2} \\
\texttt{║◌◌├──┤T║ 3} \\
\texttt{╟──┘M◌├─╢ 4} \\
\texttt{║O◌◌◌◌│S║ 5} \\
\texttt{╟─┐◌┌─┼─╢ 6} \\
\texttt{║◌│◌│P│◌║ 7} \\
\texttt{╚═╧═╧═╧═╝} \\
\texttt{```} \\
\texttt{Obecnie grasz rundę 9 z maksymalnie 28.} \\
\texttt{Łącznie zgromadziliście 7 z 16 kar.} \\
\texttt{Drugi gracz przesunął obiekt na swojej siatce.} \\
\\ 
\\ 
\texttt{Jakie jest twoje następne polecenie?} \\
            }
        }
    }
    & \\ \\

    \theutterance \stepcounter{utterance}  
    & \multicolumn{4}{p{0.6\linewidth}}{
        \cellcolor[rgb]{0.9,0.9,0.9}{
            \makecell[{{p{\linewidth}}}]{
                \texttt{\tiny{[P1$\rangle$GM]}}
                \texttt{POWIEDZ: U mnie kolumna 7 w wierszach 2{-}6 jest w całości zajęta przez krawędzie, więc P nie może wejść na (6,7). Zostawię go na (5,7) jako kompromis. Moje pozostałe cele: A=(6,1), T=(2,1), M=(5,1), S=(7,1), O=(1,7). Przesuwam teraz O z (1,5) do (1,7).} \\
            }
        }
    }
    & & \\ \\

    \theutterance \stepcounter{utterance}  
    & & & \multicolumn{2}{p{0.3\linewidth}}{
        \cellcolor[rgb]{0.9,0.9,0.9}{
            \makecell[{{p{\linewidth}}}]{
                \texttt{\tiny{[GM$|$GM]}}
                \texttt{POWIEDZ: U mnie kolumna 7 w wierszach 2{-}6 jest w całości zajęta przez krawędzie, więc P nie może wejść na (6,7). Zostawię go na (5,7) jako kompromis. Moje pozostałe cele: A=(6,1), T=(2,1), M=(5,1), S=(7,1), O=(1,7). Przesuwam teraz O z (1,5) do (1,7).} \\
            }
        }
    }
    & & \\ \\

    \theutterance \stepcounter{utterance}  
    & & \multicolumn{4}{p{0.6\linewidth}}{
        \cellcolor[rgb]{0.9,0.9,0.9}{
            \makecell[{{p{\linewidth}}}]{
                \texttt{\tiny{[GM$\rangle$P2]}}
                \texttt{Pomyślnie przeniesiono 'S' do (7,1). Twoja zaktualizowana siatka wygląda następująco:} \\
\texttt{```} \\
\texttt{ 1234567} \\
\texttt{╔══╤════╗} \\
\texttt{║◌T│◌MAS║ 1} \\
\texttt{╟──┤◌◌┌─╢ 2} \\
\texttt{║◌◌├──┤D║ 3} \\
\texttt{╟──┘◌◌├─╢ 4} \\
\texttt{║◌P◌◌◌│◌║ 5} \\
\texttt{╟─┐◌┌─┼─╢ 6} \\
\texttt{║O│◌│◌│◌║ 7} \\
\texttt{╚═╧═╧═╧═╝} \\
\texttt{```} \\
\texttt{Obecnie grasz rundę 10 z maksymalnie 28.} \\
\texttt{Łącznie zgromadziliście 7 z 16 kar.} \\
\texttt{Drugi gracz wysłał tę wiadomość:} \\
\\ 
\\ 
\texttt{"U mnie kolumna 7 w wierszach 2{-}6 jest w całości zajęta przez krawędzie, więc P nie może wejść na (6,7). Zostawię go na (5,7) jako kompromis. Moje pozostałe cele: A=(6,1), T=(2,1), M=(5,1), S=(7,1), O=(1,7). Przesuwam teraz O z (1,5) do (1,7)."} \\
\\ 
\\ 
\texttt{Jakie jest twoje następne polecenie?} \\
            }
        }
    }
    & \\ \\

    \theutterance \stepcounter{utterance}  
    & & & \multicolumn{4}{p{0.6\linewidth}}{
        \cellcolor[rgb]{0.9,0.9,0.9}{
            \makecell[{{p{\linewidth}}}]{
                \texttt{\tiny{[GM$\langle$P2]}}
                \texttt{POWIEDZ: Akceptuję P=(5,7). Przesunę u siebie P na (5,7). D zostawiam na (7,3), chyba że masz wolne (5,2); jeśli nie, zostaw u siebie D=(7,3)?} \\
            }
        }
    }
     \\ \\

    \theutterance \stepcounter{utterance}  
    & & & \multicolumn{2}{p{0.3\linewidth}}{
        \cellcolor[rgb]{0.9,0.9,0.9}{
            \makecell[{{p{\linewidth}}}]{
                \texttt{\tiny{[GM$|$GM]}}
                \texttt{POWIEDZ: Akceptuję P=(5,7). Przesunę u siebie P na (5,7). D zostawiam na (7,3), chyba że masz wolne (5,2); jeśli nie, zostaw u siebie D=(7,3)?} \\
            }
        }
    }
    & & \\ \\

    \theutterance \stepcounter{utterance}  
    & & \multicolumn{4}{p{0.6\linewidth}}{
        \cellcolor[rgb]{0.9,0.9,0.9}{
            \makecell[{{p{\linewidth}}}]{
                \texttt{\tiny{[P1$\langle$GM]}}
                \texttt{Twoja wiadomość została przekazana drugiemu graczowi.} \\
\\ 
\\ 
\texttt{Obecnie grasz rundę 10 z maksymalnie 28.} \\
\texttt{Łącznie zgromadziliście 7 z 16 kar.} \\
\texttt{Drugi gracz wysłał tę wiadomość:} \\
\\ 
\\ 
\texttt{"Akceptuję P=(5,7). Przesunę u siebie P na (5,7). D zostawiam na (7,3), chyba że masz wolne (5,2); jeśli nie, zostaw u siebie D=(7,3)?"} \\
\\ 
\\ 
\texttt{Jakie jest twoje następne polecenie?} \\
            }
        }
    }
    & \\ \\

    \theutterance \stepcounter{utterance}  
    & \multicolumn{4}{p{0.6\linewidth}}{
        \cellcolor[rgb]{0.9,0.9,0.9}{
            \makecell[{{p{\linewidth}}}]{
                \texttt{\tiny{[P1$\rangle$GM]}}
                \texttt{RUCH: O, (1, 7)} \\
            }
        }
    }
    & & \\ \\

    \theutterance \stepcounter{utterance}  
    & & & \multicolumn{2}{p{0.3\linewidth}}{
        \cellcolor[rgb]{0.9,0.9,0.9}{
            \makecell[{{p{\linewidth}}}]{
                \texttt{\tiny{[GM$|$GM]}}
                \texttt{RUCH: O, (1, 7)} \\
            }
        }
    }
    & & \\ \\

    \theutterance \stepcounter{utterance}  
    & & & \multicolumn{2}{p{0.3\linewidth}}{
        \cellcolor[rgb]{0.9,0.9,0.9}{
            \makecell[{{p{\linewidth}}}]{
                \texttt{\tiny{[GM$|$GM]}}
                \texttt{O moved to (1, 7)} \\
            }
        }
    }
    & & \\ \\

    \theutterance \stepcounter{utterance}  
    & & \multicolumn{4}{p{0.6\linewidth}}{
        \cellcolor[rgb]{0.9,0.9,0.9}{
            \makecell[{{p{\linewidth}}}]{
                \texttt{\tiny{[GM$\rangle$P2]}}
                \texttt{Twoja wiadomość została przekazana drugiemu graczowi.} \\
\\ 
\\ 
\texttt{Obecnie grasz rundę 11 z maksymalnie 28.} \\
\texttt{Łącznie zgromadziliście 7 z 16 kar.} \\
\texttt{Drugi gracz przesunął obiekt na swojej siatce.} \\
\\ 
\\ 
\texttt{Jakie jest twoje następne polecenie?} \\
            }
        }
    }
    & \\ \\

    \theutterance \stepcounter{utterance}  
    & & & \multicolumn{4}{p{0.6\linewidth}}{
        \cellcolor[rgb]{0.9,0.9,0.9}{
            \makecell[{{p{\linewidth}}}]{
                \texttt{\tiny{[GM$\langle$P2]}}
                \texttt{RUCH: P, (5,7)} \\
            }
        }
    }
     \\ \\

    \theutterance \stepcounter{utterance}  
    & & & \multicolumn{2}{p{0.3\linewidth}}{
        \cellcolor[rgb]{0.9,0.9,0.9}{
            \makecell[{{p{\linewidth}}}]{
                \texttt{\tiny{[GM$|$GM]}}
                \texttt{RUCH: P, (5,7)} \\
            }
        }
    }
    & & \\ \\

    \theutterance \stepcounter{utterance}  
    & & & \multicolumn{2}{p{0.3\linewidth}}{
        \cellcolor[rgb]{0.9,0.9,0.9}{
            \makecell[{{p{\linewidth}}}]{
                \texttt{\tiny{[GM$|$GM]}}
                \texttt{P moved to (5, 7)} \\
            }
        }
    }
    & & \\ \\

    \theutterance \stepcounter{utterance}  
    & & \multicolumn{4}{p{0.6\linewidth}}{
        \cellcolor[rgb]{0.9,0.9,0.9}{
            \makecell[{{p{\linewidth}}}]{
                \texttt{\tiny{[P1$\langle$GM]}}
                \texttt{Pomyślnie przeniesiono 'O' do (1,7). Twoja zaktualizowana siatka wygląda następująco:} \\
\texttt{```} \\
\texttt{ 1234567} \\
\texttt{╔══╤════╗} \\
\texttt{║◌◌│◌◌◌A║ 1} \\
\texttt{╟──┤◌D┌─╢ 2} \\
\texttt{║◌◌├──┤T║ 3} \\
\texttt{╟──┘M◌├─╢ 4} \\
\texttt{║◌◌◌◌◌│S║ 5} \\
\texttt{╟─┐◌┌─┼─╢ 6} \\
\texttt{║O│◌│P│◌║ 7} \\
\texttt{╚═╧═╧═╧═╝} \\
\texttt{```} \\
\texttt{Obecnie grasz rundę 11 z maksymalnie 28.} \\
\texttt{Łącznie zgromadziliście 7 z 16 kar.} \\
\texttt{Drugi gracz przesunął obiekt na swojej siatce.} \\
\\ 
\\ 
\texttt{Jakie jest twoje następne polecenie?} \\
            }
        }
    }
    & \\ \\

    \theutterance \stepcounter{utterance}  
    & \multicolumn{4}{p{0.6\linewidth}}{
        \cellcolor[rgb]{0.9,0.9,0.9}{
            \makecell[{{p{\linewidth}}}]{
                \texttt{\tiny{[P1$\rangle$GM]}}
                \texttt{POWIEDZ: Moje D jest już na (5,2) i to pole jest u mnie puste, więc zostawiam go tam. Przesuwam teraz T z (7,3) do (3,3).} \\
            }
        }
    }
    & & \\ \\

    \theutterance \stepcounter{utterance}  
    & & & \multicolumn{2}{p{0.3\linewidth}}{
        \cellcolor[rgb]{0.9,0.9,0.9}{
            \makecell[{{p{\linewidth}}}]{
                \texttt{\tiny{[GM$|$GM]}}
                \texttt{POWIEDZ: Moje D jest już na (5,2) i to pole jest u mnie puste, więc zostawiam go tam. Przesuwam teraz T z (7,3) do (3,3).} \\
            }
        }
    }
    & & \\ \\

    \theutterance \stepcounter{utterance}  
    & & \multicolumn{4}{p{0.6\linewidth}}{
        \cellcolor[rgb]{0.9,0.9,0.9}{
            \makecell[{{p{\linewidth}}}]{
                \texttt{\tiny{[GM$\rangle$P2]}}
                \texttt{Pomyślnie przeniesiono 'P' do (5,7). Twoja zaktualizowana siatka wygląda następująco:} \\
\texttt{```} \\
\texttt{ 1234567} \\
\texttt{╔══╤════╗} \\
\texttt{║◌T│◌MAS║ 1} \\
\texttt{╟──┤◌◌┌─╢ 2} \\
\texttt{║◌◌├──┤D║ 3} \\
\texttt{╟──┘◌◌├─╢ 4} \\
\texttt{║◌◌◌◌◌│◌║ 5} \\
\texttt{╟─┐◌┌─┼─╢ 6} \\
\texttt{║O│◌│P│◌║ 7} \\
\texttt{╚═╧═╧═╧═╝} \\
\texttt{```} \\
\texttt{Obecnie grasz rundę 12 z maksymalnie 28.} \\
\texttt{Łącznie zgromadziliście 7 z 16 kar.} \\
\texttt{Drugi gracz wysłał tę wiadomość:} \\
\\ 
\\ 
\texttt{"Moje D jest już na (5,2) i to pole jest u mnie puste, więc zostawiam go tam. Przesuwam teraz T z (7,3) do (3,3)."} \\
\\ 
\\ 
\texttt{Jakie jest twoje następne polecenie?} \\
            }
        }
    }
    & \\ \\

    \theutterance \stepcounter{utterance}  
    & & & \multicolumn{4}{p{0.6\linewidth}}{
        \cellcolor[rgb]{0.9,0.9,0.9}{
            \makecell[{{p{\linewidth}}}]{
                \texttt{\tiny{[GM$\langle$P2]}}
                \texttt{POWIEDZ: Nie mogę przesunąć D na (5,2), bo u mnie to pole zawiera krawędź. Najbliższe wolne pole to (1,5). Czy postawisz u siebie D na (1,5)? Wtedy zakończymy równo.} \\
            }
        }
    }
     \\ \\

    \theutterance \stepcounter{utterance}  
    & & & \multicolumn{2}{p{0.3\linewidth}}{
        \cellcolor[rgb]{0.9,0.9,0.9}{
            \makecell[{{p{\linewidth}}}]{
                \texttt{\tiny{[GM$|$GM]}}
                \texttt{POWIEDZ: Nie mogę przesunąć D na (5,2), bo u mnie to pole zawiera krawędź. Najbliższe wolne pole to (1,5). Czy postawisz u siebie D na (1,5)? Wtedy zakończymy równo.} \\
            }
        }
    }
    & & \\ \\

    \theutterance \stepcounter{utterance}  
    & & \multicolumn{4}{p{0.6\linewidth}}{
        \cellcolor[rgb]{0.9,0.9,0.9}{
            \makecell[{{p{\linewidth}}}]{
                \texttt{\tiny{[P1$\langle$GM]}}
                \texttt{Twoja wiadomość została przekazana drugiemu graczowi.} \\
\\ 
\\ 
\texttt{Obecnie grasz rundę 12 z maksymalnie 28.} \\
\texttt{Łącznie zgromadziliście 7 z 16 kar.} \\
\texttt{Drugi gracz wysłał tę wiadomość:} \\
\\ 
\\ 
\texttt{"Nie mogę przesunąć D na (5,2), bo u mnie to pole zawiera krawędź. Najbliższe wolne pole to (1,5). Czy postawisz u siebie D na (1,5)? Wtedy zakończymy równo."} \\
\\ 
\\ 
\texttt{Jakie jest twoje następne polecenie?} \\
            }
        }
    }
    & \\ \\

    \theutterance \stepcounter{utterance}  
    & \multicolumn{4}{p{0.6\linewidth}}{
        \cellcolor[rgb]{0.9,0.9,0.9}{
            \makecell[{{p{\linewidth}}}]{
                \texttt{\tiny{[P1$\rangle$GM]}}
                \texttt{RUCH: D, (1, 5)} \\
            }
        }
    }
    & & \\ \\

    \theutterance \stepcounter{utterance}  
    & & & \multicolumn{2}{p{0.3\linewidth}}{
        \cellcolor[rgb]{0.9,0.9,0.9}{
            \makecell[{{p{\linewidth}}}]{
                \texttt{\tiny{[GM$|$GM]}}
                \texttt{RUCH: D, (1, 5)} \\
            }
        }
    }
    & & \\ \\

    \theutterance \stepcounter{utterance}  
    & & & \multicolumn{2}{p{0.3\linewidth}}{
        \cellcolor[rgb]{0.9,0.9,0.9}{
            \makecell[{{p{\linewidth}}}]{
                \texttt{\tiny{[GM$|$GM]}}
                \texttt{D moved to (1, 5)} \\
            }
        }
    }
    & & \\ \\

    \theutterance \stepcounter{utterance}  
    & & \multicolumn{4}{p{0.6\linewidth}}{
        \cellcolor[rgb]{0.9,0.9,0.9}{
            \makecell[{{p{\linewidth}}}]{
                \texttt{\tiny{[GM$\rangle$P2]}}
                \texttt{Twoja wiadomość została przekazana drugiemu graczowi.} \\
\\ 
\\ 
\texttt{Obecnie grasz rundę 13 z maksymalnie 28.} \\
\texttt{Łącznie zgromadziliście 7 z 16 kar.} \\
\texttt{Drugi gracz przesunął obiekt na swojej siatce.} \\
\\ 
\\ 
\texttt{Jakie jest twoje następne polecenie?} \\
            }
        }
    }
    & \\ \\

    \theutterance \stepcounter{utterance}  
    & & & \multicolumn{4}{p{0.6\linewidth}}{
        \cellcolor[rgb]{0.9,0.9,0.9}{
            \makecell[{{p{\linewidth}}}]{
                \texttt{\tiny{[GM$\langle$P2]}}
                \texttt{RUCH: D, (1,5)} \\
            }
        }
    }
     \\ \\

    \theutterance \stepcounter{utterance}  
    & & & \multicolumn{2}{p{0.3\linewidth}}{
        \cellcolor[rgb]{0.9,0.9,0.9}{
            \makecell[{{p{\linewidth}}}]{
                \texttt{\tiny{[GM$|$GM]}}
                \texttt{RUCH: D, (1,5)} \\
            }
        }
    }
    & & \\ \\

    \theutterance \stepcounter{utterance}  
    & & & \multicolumn{2}{p{0.3\linewidth}}{
        \cellcolor[rgb]{0.9,0.9,0.9}{
            \makecell[{{p{\linewidth}}}]{
                \texttt{\tiny{[GM$|$GM]}}
                \texttt{D moved to (1, 5)} \\
            }
        }
    }
    & & \\ \\

    \theutterance \stepcounter{utterance}  
    & & \multicolumn{4}{p{0.6\linewidth}}{
        \cellcolor[rgb]{0.9,0.9,0.9}{
            \makecell[{{p{\linewidth}}}]{
                \texttt{\tiny{[P1$\langle$GM]}}
                \texttt{Pomyślnie przeniesiono 'D' do (1,5). Twoja zaktualizowana siatka wygląda następująco:} \\
\texttt{```} \\
\texttt{ 1234567} \\
\texttt{╔══╤════╗} \\
\texttt{║◌◌│◌◌◌A║ 1} \\
\texttt{╟──┤◌◌┌─╢ 2} \\
\texttt{║◌◌├──┤T║ 3} \\
\texttt{╟──┘M◌├─╢ 4} \\
\texttt{║D◌◌◌◌│S║ 5} \\
\texttt{╟─┐◌┌─┼─╢ 6} \\
\texttt{║O│◌│P│◌║ 7} \\
\texttt{╚═╧═╧═╧═╝} \\
\texttt{```} \\
\texttt{Obecnie grasz rundę 13 z maksymalnie 28.} \\
\texttt{Łącznie zgromadziliście 7 z 16 kar.} \\
\texttt{Drugi gracz przesunął obiekt na swojej siatce.} \\
\\ 
\\ 
\texttt{Jakie jest twoje następne polecenie?} \\
            }
        }
    }
    & \\ \\

    \theutterance \stepcounter{utterance}  
    & \multicolumn{4}{p{0.6\linewidth}}{
        \cellcolor[rgb]{0.9,0.9,0.9}{
            \makecell[{{p{\linewidth}}}]{
                \texttt{\tiny{[P1$\rangle$GM]}}
                \texttt{RUCH: T, (6, 3)} \\
            }
        }
    }
    & & \\ \\

    \theutterance \stepcounter{utterance}  
    & & & \multicolumn{2}{p{0.3\linewidth}}{
        \cellcolor[rgb]{0.9,0.9,0.9}{
            \makecell[{{p{\linewidth}}}]{
                \texttt{\tiny{[GM$|$GM]}}
                \texttt{RUCH: T, (6, 3)} \\
            }
        }
    }
    & & \\ \\

    \theutterance \stepcounter{utterance}  
    & & & \multicolumn{2}{p{0.3\linewidth}}{
        \cellcolor[rgb]{0.9,0.9,0.9}{
            \makecell[{{p{\linewidth}}}]{
                \texttt{\tiny{[GM$|$GM]}}
                \texttt{Kara: (6,3) nie jest puste, ale zawiera '┤'.} \\
\\ 
\texttt{Łącznie zgromadziliście 8 z 16 kar.} \\
\texttt{Spróbuj ponownie!} \\
            }
        }
    }
    & & \\ \\

    \theutterance \stepcounter{utterance}  
    & & \multicolumn{4}{p{0.6\linewidth}}{
        \cellcolor[rgb]{0.9,0.9,0.9}{
            \makecell[{{p{\linewidth}}}]{
                \texttt{\tiny{[P1$\langle$GM]}}
                \texttt{Kara: (6,3) nie jest puste, ale zawiera '┤'.} \\
\\ 
\texttt{Łącznie zgromadziliście 8 z 16 kar.} \\
\texttt{Spróbuj ponownie!} \\
            }
        }
    }
    & \\ \\

    \theutterance \stepcounter{utterance}  
    & \multicolumn{4}{p{0.6\linewidth}}{
        \cellcolor[rgb]{0.9,0.9,0.9}{
            \makecell[{{p{\linewidth}}}]{
                \texttt{\tiny{[P1$\rangle$GM]}}
                \texttt{RUCH: T, (3, 3)} \\
            }
        }
    }
    & & \\ \\

    \theutterance \stepcounter{utterance}  
    & & & \multicolumn{2}{p{0.3\linewidth}}{
        \cellcolor[rgb]{0.9,0.9,0.9}{
            \makecell[{{p{\linewidth}}}]{
                \texttt{\tiny{[GM$|$GM]}}
                \texttt{RUCH: T, (3, 3)} \\
            }
        }
    }
    & & \\ \\

    \theutterance \stepcounter{utterance}  
    & & & \multicolumn{2}{p{0.3\linewidth}}{
        \cellcolor[rgb]{0.9,0.9,0.9}{
            \makecell[{{p{\linewidth}}}]{
                \texttt{\tiny{[GM$|$GM]}}
                \texttt{Kara: (3,3) nie jest puste, ale zawiera '├'.} \\
\\ 
\texttt{Łącznie zgromadziliście 9 z 16 kar.} \\
\texttt{Spróbuj ponownie!} \\
            }
        }
    }
    & & \\ \\

    \theutterance \stepcounter{utterance}  
    & & \multicolumn{4}{p{0.6\linewidth}}{
        \cellcolor[rgb]{0.9,0.9,0.9}{
            \makecell[{{p{\linewidth}}}]{
                \texttt{\tiny{[P1$\langle$GM]}}
                \texttt{Kara: (3,3) nie jest puste, ale zawiera '├'.} \\
\\ 
\texttt{Łącznie zgromadziliście 9 z 16 kar.} \\
\texttt{Spróbuj ponownie!} \\
            }
        }
    }
    & \\ \\

    \theutterance \stepcounter{utterance}  
    & \multicolumn{4}{p{0.6\linewidth}}{
        \cellcolor[rgb]{0.9,0.9,0.9}{
            \makecell[{{p{\linewidth}}}]{
                \texttt{\tiny{[P1$\rangle$GM]}}
                \texttt{POWIEDZ: U mnie T=(2,1) zawiera górną krawędź, a wiersz 3 jest pełen ścian. Zostawię T na (7,3) jako kompromis. Przesuwam teraz M z (2,4) do (5,1).} \\
            }
        }
    }
    & & \\ \\

    \theutterance \stepcounter{utterance}  
    & & & \multicolumn{2}{p{0.3\linewidth}}{
        \cellcolor[rgb]{0.9,0.9,0.9}{
            \makecell[{{p{\linewidth}}}]{
                \texttt{\tiny{[GM$|$GM]}}
                \texttt{POWIEDZ: U mnie T=(2,1) zawiera górną krawędź, a wiersz 3 jest pełen ścian. Zostawię T na (7,3) jako kompromis. Przesuwam teraz M z (2,4) do (5,1).} \\
            }
        }
    }
    & & \\ \\

    \theutterance \stepcounter{utterance}  
    & & \multicolumn{4}{p{0.6\linewidth}}{
        \cellcolor[rgb]{0.9,0.9,0.9}{
            \makecell[{{p{\linewidth}}}]{
                \texttt{\tiny{[GM$\rangle$P2]}}
                \texttt{Pomyślnie przeniesiono 'D' do (1,5). Twoja zaktualizowana siatka wygląda następująco:} \\
\texttt{```} \\
\texttt{ 1234567} \\
\texttt{╔══╤════╗} \\
\texttt{║◌T│◌MAS║ 1} \\
\texttt{╟──┤◌◌┌─╢ 2} \\
\texttt{║◌◌├──┤◌║ 3} \\
\texttt{╟──┘◌◌├─╢ 4} \\
\texttt{║D◌◌◌◌│◌║ 5} \\
\texttt{╟─┐◌┌─┼─╢ 6} \\
\texttt{║O│◌│P│◌║ 7} \\
\texttt{╚═╧═╧═╧═╝} \\
\texttt{```} \\
\texttt{Obecnie grasz rundę 14 z maksymalnie 28.} \\
\texttt{Łącznie zgromadziliście 9 z 16 kar.} \\
\texttt{Drugi gracz wysłał tę wiadomość:} \\
\\ 
\\ 
\texttt{"U mnie T=(2,1) zawiera górną krawędź, a wiersz 3 jest pełen ścian. Zostawię T na (7,3) jako kompromis. Przesuwam teraz M z (2,4) do (5,1)."} \\
\\ 
\\ 
\texttt{Jakie jest twoje następne polecenie?} \\
            }
        }
    }
    & \\ \\

    \theutterance \stepcounter{utterance}  
    & & & \multicolumn{4}{p{0.6\linewidth}}{
        \cellcolor[rgb]{0.9,0.9,0.9}{
            \makecell[{{p{\linewidth}}}]{
                \texttt{\tiny{[GM$\langle$P2]}}
                \texttt{POWIEDZ: Akceptuję T=(7,3). Zostawię u siebie T na (2,1), a D przesunąłem na (1,5). Proszę, ustaw D u siebie na (1,5). Moja finalna konfiguracja: A=(6,1), T=(2,1), M=(5,1), S=(7,1), O=(1,7), P=(5,7), D=(1,5).} \\
            }
        }
    }
     \\ \\

    \theutterance \stepcounter{utterance}  
    & & & \multicolumn{2}{p{0.3\linewidth}}{
        \cellcolor[rgb]{0.9,0.9,0.9}{
            \makecell[{{p{\linewidth}}}]{
                \texttt{\tiny{[GM$|$GM]}}
                \texttt{POWIEDZ: Akceptuję T=(7,3). Zostawię u siebie T na (2,1), a D przesunąłem na (1,5). Proszę, ustaw D u siebie na (1,5). Moja finalna konfiguracja: A=(6,1), T=(2,1), M=(5,1), S=(7,1), O=(1,7), P=(5,7), D=(1,5).} \\
            }
        }
    }
    & & \\ \\

    \theutterance \stepcounter{utterance}  
    & & \multicolumn{4}{p{0.6\linewidth}}{
        \cellcolor[rgb]{0.9,0.9,0.9}{
            \makecell[{{p{\linewidth}}}]{
                \texttt{\tiny{[P1$\langle$GM]}}
                \texttt{Twoja wiadomość została przekazana drugiemu graczowi.} \\
\\ 
\\ 
\texttt{Obecnie grasz rundę 14 z maksymalnie 28.} \\
\texttt{Łącznie zgromadziliście 9 z 16 kar.} \\
\texttt{Drugi gracz wysłał tę wiadomość:} \\
\\ 
\\ 
\texttt{"Akceptuję T=(7,3). Zostawię u siebie T na (2,1), a D przesunąłem na (1,5). Proszę, ustaw D u siebie na (1,5). Moja finalna konfiguracja: A=(6,1), T=(2,1), M=(5,1), S=(7,1), O=(1,7), P=(5,7), D=(1,5)."} \\
\\ 
\\ 
\texttt{Jakie jest twoje następne polecenie?} \\
            }
        }
    }
    & \\ \\

    \theutterance \stepcounter{utterance}  
    & \multicolumn{4}{p{0.6\linewidth}}{
        \cellcolor[rgb]{0.9,0.9,0.9}{
            \makecell[{{p{\linewidth}}}]{
                \texttt{\tiny{[P1$\rangle$GM]}}
                \texttt{RUCH: M, (5, 4)} \\
            }
        }
    }
    & & \\ \\

    \theutterance \stepcounter{utterance}  
    & & & \multicolumn{2}{p{0.3\linewidth}}{
        \cellcolor[rgb]{0.9,0.9,0.9}{
            \makecell[{{p{\linewidth}}}]{
                \texttt{\tiny{[GM$|$GM]}}
                \texttt{RUCH: M, (5, 4)} \\
            }
        }
    }
    & & \\ \\

    \theutterance \stepcounter{utterance}  
    & & & \multicolumn{2}{p{0.3\linewidth}}{
        \cellcolor[rgb]{0.9,0.9,0.9}{
            \makecell[{{p{\linewidth}}}]{
                \texttt{\tiny{[GM$|$GM]}}
                \texttt{M moved to (5, 4)} \\
            }
        }
    }
    & & \\ \\

    \theutterance \stepcounter{utterance}  
    & & \multicolumn{4}{p{0.6\linewidth}}{
        \cellcolor[rgb]{0.9,0.9,0.9}{
            \makecell[{{p{\linewidth}}}]{
                \texttt{\tiny{[GM$\rangle$P2]}}
                \texttt{Twoja wiadomość została przekazana drugiemu graczowi.} \\
\\ 
\\ 
\texttt{Obecnie grasz rundę 15 z maksymalnie 28.} \\
\texttt{Łącznie zgromadziliście 9 z 16 kar.} \\
\texttt{Drugi gracz przesunął obiekt na swojej siatce.} \\
\\ 
\\ 
\texttt{Jakie jest twoje następne polecenie?} \\
            }
        }
    }
    & \\ \\

    \theutterance \stepcounter{utterance}  
    & & & \multicolumn{4}{p{0.6\linewidth}}{
        \cellcolor[rgb]{0.9,0.9,0.9}{
            \makecell[{{p{\linewidth}}}]{
                \texttt{\tiny{[GM$\langle$P2]}}
                \texttt{RUCH: T, (7,3)} \\
            }
        }
    }
     \\ \\

    \theutterance \stepcounter{utterance}  
    & & & \multicolumn{2}{p{0.3\linewidth}}{
        \cellcolor[rgb]{0.9,0.9,0.9}{
            \makecell[{{p{\linewidth}}}]{
                \texttt{\tiny{[GM$|$GM]}}
                \texttt{RUCH: T, (7,3)} \\
            }
        }
    }
    & & \\ \\

    \theutterance \stepcounter{utterance}  
    & & & \multicolumn{2}{p{0.3\linewidth}}{
        \cellcolor[rgb]{0.9,0.9,0.9}{
            \makecell[{{p{\linewidth}}}]{
                \texttt{\tiny{[GM$|$GM]}}
                \texttt{T moved to (7, 3)} \\
            }
        }
    }
    & & \\ \\

    \theutterance \stepcounter{utterance}  
    & & \multicolumn{4}{p{0.6\linewidth}}{
        \cellcolor[rgb]{0.9,0.9,0.9}{
            \makecell[{{p{\linewidth}}}]{
                \texttt{\tiny{[P1$\langle$GM]}}
                \texttt{Pomyślnie przeniesiono 'M' do (5,4). Twoja zaktualizowana siatka wygląda następująco:} \\
\texttt{```} \\
\texttt{ 1234567} \\
\texttt{╔══╤════╗} \\
\texttt{║◌◌│◌◌◌A║ 1} \\
\texttt{╟──┤◌◌┌─╢ 2} \\
\texttt{║◌◌├──┤T║ 3} \\
\texttt{╟──┘◌M├─╢ 4} \\
\texttt{║D◌◌◌◌│S║ 5} \\
\texttt{╟─┐◌┌─┼─╢ 6} \\
\texttt{║O│◌│P│◌║ 7} \\
\texttt{╚═╧═╧═╧═╝} \\
\texttt{```} \\
\texttt{Obecnie grasz rundę 15 z maksymalnie 28.} \\
\texttt{Łącznie zgromadziliście 9 z 16 kar.} \\
\texttt{Drugi gracz przesunął obiekt na swojej siatce.} \\
\\ 
\\ 
\texttt{Jakie jest twoje następne polecenie?} \\
            }
        }
    }
    & \\ \\

    \theutterance \stepcounter{utterance}  
    & \multicolumn{4}{p{0.6\linewidth}}{
        \cellcolor[rgb]{0.9,0.9,0.9}{
            \makecell[{{p{\linewidth}}}]{
                \texttt{\tiny{[P1$\rangle$GM]}}
                \texttt{RUCH: M, (5, 1)} \\
            }
        }
    }
    & & \\ \\

    \theutterance \stepcounter{utterance}  
    & & & \multicolumn{2}{p{0.3\linewidth}}{
        \cellcolor[rgb]{0.9,0.9,0.9}{
            \makecell[{{p{\linewidth}}}]{
                \texttt{\tiny{[GM$|$GM]}}
                \texttt{RUCH: M, (5, 1)} \\
            }
        }
    }
    & & \\ \\

    \theutterance \stepcounter{utterance}  
    & & & \multicolumn{2}{p{0.3\linewidth}}{
        \cellcolor[rgb]{0.9,0.9,0.9}{
            \makecell[{{p{\linewidth}}}]{
                \texttt{\tiny{[GM$|$GM]}}
                \texttt{M moved to (5, 1)} \\
            }
        }
    }
    & & \\ \\

    \theutterance \stepcounter{utterance}  
    & & \multicolumn{4}{p{0.6\linewidth}}{
        \cellcolor[rgb]{0.9,0.9,0.9}{
            \makecell[{{p{\linewidth}}}]{
                \texttt{\tiny{[GM$\rangle$P2]}}
                \texttt{Pomyślnie przeniesiono 'T' do (7,3). Twoja zaktualizowana siatka wygląda następująco:} \\
\texttt{```} \\
\texttt{ 1234567} \\
\texttt{╔══╤════╗} \\
\texttt{║◌◌│◌MAS║ 1} \\
\texttt{╟──┤◌◌┌─╢ 2} \\
\texttt{║◌◌├──┤T║ 3} \\
\texttt{╟──┘◌◌├─╢ 4} \\
\texttt{║D◌◌◌◌│◌║ 5} \\
\texttt{╟─┐◌┌─┼─╢ 6} \\
\texttt{║O│◌│P│◌║ 7} \\
\texttt{╚═╧═╧═╧═╝} \\
\texttt{```} \\
\texttt{Obecnie grasz rundę 16 z maksymalnie 28.} \\
\texttt{Łącznie zgromadziliście 9 z 16 kar.} \\
\texttt{Drugi gracz przesunął obiekt na swojej siatce.} \\
\\ 
\\ 
\texttt{Jakie jest twoje następne polecenie?} \\
            }
        }
    }
    & \\ \\

    \theutterance \stepcounter{utterance}  
    & & & \multicolumn{4}{p{0.6\linewidth}}{
        \cellcolor[rgb]{0.9,0.9,0.9}{
            \makecell[{{p{\linewidth}}}]{
                \texttt{\tiny{[GM$\langle$P2]}}
                \texttt{POWIEDZ: T=(7,3) gotowe. Czekam, aż ustawisz D na (1,5), a potem proszę o zakończenie.} \\
            }
        }
    }
     \\ \\

    \theutterance \stepcounter{utterance}  
    & & & \multicolumn{2}{p{0.3\linewidth}}{
        \cellcolor[rgb]{0.9,0.9,0.9}{
            \makecell[{{p{\linewidth}}}]{
                \texttt{\tiny{[GM$|$GM]}}
                \texttt{POWIEDZ: T=(7,3) gotowe. Czekam, aż ustawisz D na (1,5), a potem proszę o zakończenie.} \\
            }
        }
    }
    & & \\ \\

    \theutterance \stepcounter{utterance}  
    & & \multicolumn{4}{p{0.6\linewidth}}{
        \cellcolor[rgb]{0.9,0.9,0.9}{
            \makecell[{{p{\linewidth}}}]{
                \texttt{\tiny{[P1$\langle$GM]}}
                \texttt{Pomyślnie przeniesiono 'M' do (5,1). Twoja zaktualizowana siatka wygląda następująco:} \\
\texttt{```} \\
\texttt{ 1234567} \\
\texttt{╔══╤════╗} \\
\texttt{║◌◌│◌M◌A║ 1} \\
\texttt{╟──┤◌◌┌─╢ 2} \\
\texttt{║◌◌├──┤T║ 3} \\
\texttt{╟──┘◌◌├─╢ 4} \\
\texttt{║D◌◌◌◌│S║ 5} \\
\texttt{╟─┐◌┌─┼─╢ 6} \\
\texttt{║O│◌│P│◌║ 7} \\
\texttt{╚═╧═╧═╧═╝} \\
\texttt{```} \\
\texttt{Obecnie grasz rundę 16 z maksymalnie 28.} \\
\texttt{Łącznie zgromadziliście 9 z 16 kar.} \\
\texttt{Drugi gracz wysłał tę wiadomość:} \\
\\ 
\\ 
\texttt{"T=(7,3) gotowe. Czekam, aż ustawisz D na (1,5), a potem proszę o zakończenie."} \\
\\ 
\\ 
\texttt{Jakie jest twoje następne polecenie?} \\
            }
        }
    }
    & \\ \\

    \theutterance \stepcounter{utterance}  
    & \multicolumn{4}{p{0.6\linewidth}}{
        \cellcolor[rgb]{0.9,0.9,0.9}{
            \makecell[{{p{\linewidth}}}]{
                \texttt{\tiny{[P1$\rangle$GM]}}
                \texttt{POWIEDZ: D jest już na (1,5), a T na (7,3). Brakuje tylko przesunięcia A z (7,1) do (6,1). Potem poproszę o zakończenie.} \\
            }
        }
    }
    & & \\ \\

    \theutterance \stepcounter{utterance}  
    & & & \multicolumn{2}{p{0.3\linewidth}}{
        \cellcolor[rgb]{0.9,0.9,0.9}{
            \makecell[{{p{\linewidth}}}]{
                \texttt{\tiny{[GM$|$GM]}}
                \texttt{POWIEDZ: D jest już na (1,5), a T na (7,3). Brakuje tylko przesunięcia A z (7,1) do (6,1). Potem poproszę o zakończenie.} \\
            }
        }
    }
    & & \\ \\

    \theutterance \stepcounter{utterance}  
    & & \multicolumn{4}{p{0.6\linewidth}}{
        \cellcolor[rgb]{0.9,0.9,0.9}{
            \makecell[{{p{\linewidth}}}]{
                \texttt{\tiny{[GM$\rangle$P2]}}
                \texttt{Twoja wiadomość została przekazana drugiemu graczowi.} \\
\\ 
\\ 
\texttt{Obecnie grasz rundę 17 z maksymalnie 28.} \\
\texttt{Łącznie zgromadziliście 9 z 16 kar.} \\
\texttt{Drugi gracz wysłał tę wiadomość:} \\
\\ 
\\ 
\texttt{"D jest już na (1,5), a T na (7,3). Brakuje tylko przesunięcia A z (7,1) do (6,1). Potem poproszę o zakończenie."} \\
\\ 
\\ 
\texttt{Jakie jest twoje następne polecenie?} \\
            }
        }
    }
    & \\ \\

    \theutterance \stepcounter{utterance}  
    & & & \multicolumn{4}{p{0.6\linewidth}}{
        \cellcolor[rgb]{0.9,0.9,0.9}{
            \makecell[{{p{\linewidth}}}]{
                \texttt{\tiny{[GM$\langle$P2]}}
                \texttt{POWIEDZ: U mnie A jest już na (6,1), a reszta na miejscu: T=(7,3), M=(5,1), S=(7,1), O=(1,7), P=(5,7), D=(1,5). Możemy zakończyć.} \\
            }
        }
    }
     \\ \\

    \theutterance \stepcounter{utterance}  
    & & & \multicolumn{2}{p{0.3\linewidth}}{
        \cellcolor[rgb]{0.9,0.9,0.9}{
            \makecell[{{p{\linewidth}}}]{
                \texttt{\tiny{[GM$|$GM]}}
                \texttt{POWIEDZ: U mnie A jest już na (6,1), a reszta na miejscu: T=(7,3), M=(5,1), S=(7,1), O=(1,7), P=(5,7), D=(1,5). Możemy zakończyć.} \\
            }
        }
    }
    & & \\ \\

    \theutterance \stepcounter{utterance}  
    & & \multicolumn{4}{p{0.6\linewidth}}{
        \cellcolor[rgb]{0.9,0.9,0.9}{
            \makecell[{{p{\linewidth}}}]{
                \texttt{\tiny{[P1$\langle$GM]}}
                \texttt{Twoja wiadomość została przekazana drugiemu graczowi.} \\
\\ 
\\ 
\texttt{Obecnie grasz rundę 17 z maksymalnie 28.} \\
\texttt{Łącznie zgromadziliście 9 z 16 kar.} \\
\texttt{Drugi gracz wysłał tę wiadomość:} \\
\\ 
\\ 
\texttt{"U mnie A jest już na (6,1), a reszta na miejscu: T=(7,3), M=(5,1), S=(7,1), O=(1,7), P=(5,7), D=(1,5). Możemy zakończyć."} \\
\\ 
\\ 
\texttt{Jakie jest twoje następne polecenie?} \\
            }
        }
    }
    & \\ \\

    \theutterance \stepcounter{utterance}  
    & \multicolumn{4}{p{0.6\linewidth}}{
        \cellcolor[rgb]{0.9,0.9,0.9}{
            \makecell[{{p{\linewidth}}}]{
                \texttt{\tiny{[P1$\rangle$GM]}}
                \texttt{RUCH: A, (6, 1)} \\
            }
        }
    }
    & & \\ \\

    \theutterance \stepcounter{utterance}  
    & & & \multicolumn{2}{p{0.3\linewidth}}{
        \cellcolor[rgb]{0.9,0.9,0.9}{
            \makecell[{{p{\linewidth}}}]{
                \texttt{\tiny{[GM$|$GM]}}
                \texttt{RUCH: A, (6, 1)} \\
            }
        }
    }
    & & \\ \\

    \theutterance \stepcounter{utterance}  
    & & & \multicolumn{2}{p{0.3\linewidth}}{
        \cellcolor[rgb]{0.9,0.9,0.9}{
            \makecell[{{p{\linewidth}}}]{
                \texttt{\tiny{[GM$|$GM]}}
                \texttt{A moved to (6, 1)} \\
            }
        }
    }
    & & \\ \\

    \theutterance \stepcounter{utterance}  
    & & \multicolumn{4}{p{0.6\linewidth}}{
        \cellcolor[rgb]{0.9,0.9,0.9}{
            \makecell[{{p{\linewidth}}}]{
                \texttt{\tiny{[GM$\rangle$P2]}}
                \texttt{Twoja wiadomość została przekazana drugiemu graczowi.} \\
\\ 
\\ 
\texttt{Obecnie grasz rundę 18 z maksymalnie 28.} \\
\texttt{Łącznie zgromadziliście 9 z 16 kar.} \\
\texttt{Drugi gracz przesunął obiekt na swojej siatce.} \\
\\ 
\\ 
\texttt{Jakie jest twoje następne polecenie?} \\
            }
        }
    }
    & \\ \\

    \theutterance \stepcounter{utterance}  
    & & & \multicolumn{4}{p{0.6\linewidth}}{
        \cellcolor[rgb]{0.9,0.9,0.9}{
            \makecell[{{p{\linewidth}}}]{
                \texttt{\tiny{[GM$\langle$P2]}}
                \texttt{POWIEDZ: zakończono?} \\
            }
        }
    }
     \\ \\

    \theutterance \stepcounter{utterance}  
    & & & \multicolumn{2}{p{0.3\linewidth}}{
        \cellcolor[rgb]{0.9,0.9,0.9}{
            \makecell[{{p{\linewidth}}}]{
                \texttt{\tiny{[GM$|$GM]}}
                \texttt{POWIEDZ: zakończono?} \\
            }
        }
    }
    & & \\ \\

    \theutterance \stepcounter{utterance}  
    & & \multicolumn{4}{p{0.6\linewidth}}{
        \cellcolor[rgb]{0.9,0.9,0.9}{
            \makecell[{{p{\linewidth}}}]{
                \texttt{\tiny{[P1$\langle$GM]}}
                \texttt{Pomyślnie przeniesiono 'A' do (6,1). Twoja zaktualizowana siatka wygląda następująco:} \\
\texttt{```} \\
\texttt{ 1234567} \\
\texttt{╔══╤════╗} \\
\texttt{║◌◌│◌MA◌║ 1} \\
\texttt{╟──┤◌◌┌─╢ 2} \\
\texttt{║◌◌├──┤T║ 3} \\
\texttt{╟──┘◌◌├─╢ 4} \\
\texttt{║D◌◌◌◌│S║ 5} \\
\texttt{╟─┐◌┌─┼─╢ 6} \\
\texttt{║O│◌│P│◌║ 7} \\
\texttt{╚═╧═╧═╧═╝} \\
\texttt{```} \\
\texttt{Obecnie grasz rundę 18 z maksymalnie 28.} \\
\texttt{Łącznie zgromadziliście 9 z 16 kar.} \\
\texttt{Drugi gracz wysłał tę wiadomość:} \\
\\ 
\\ 
\texttt{"zakończono?"} \\
\\ 
\\ 
\texttt{Jakie jest twoje następne polecenie?} \\
            }
        }
    }
    & \\ \\

    \theutterance \stepcounter{utterance}  
    & \multicolumn{4}{p{0.6\linewidth}}{
        \cellcolor[rgb]{0.9,0.9,0.9}{
            \makecell[{{p{\linewidth}}}]{
                \texttt{\tiny{[P1$\rangle$GM]}}
                \texttt{POWIEDZ: zakończono!} \\
            }
        }
    }
    & & \\ \\

    \theutterance \stepcounter{utterance}  
    & & & \multicolumn{2}{p{0.3\linewidth}}{
        \cellcolor[rgb]{0.9,0.9,0.9}{
            \makecell[{{p{\linewidth}}}]{
                \texttt{\tiny{[GM$|$GM]}}
                \texttt{POWIEDZ: zakończono!} \\
            }
        }
    }
    & & \\ \\

    \theutterance \stepcounter{utterance}  
    & & & \multicolumn{2}{p{0.3\linewidth}}{
        \cellcolor[rgb]{0.9,0.9,0.9}{
            \makecell[{{p{\linewidth}}}]{
                \texttt{\tiny{[GM$|$GM]}}
                \texttt{true} \\
            }
        }
    }
    & & \\ \\

    \theutterance \stepcounter{utterance}  
    & & & \multicolumn{2}{p{0.3\linewidth}}{
        \cellcolor[rgb]{0.9,0.9,0.9}{
            \makecell[{{p{\linewidth}}}]{
                \texttt{\tiny{[GM$|$GM]}}
                \texttt{Initial states:} \\
\texttt{```} \\
\texttt{Player 1:} \\
\texttt{ 1234567} \\
\texttt{╔══╤════╗} \\
\texttt{║◌◌│◌◌◌A║ 1} \\
\texttt{╟──┤◌◌┌─╢ 2} \\
\texttt{║◌◌├──┤T║ 3} \\
\texttt{╟──┘M◌├─╢ 4} \\
\texttt{║P◌OD◌│S║ 5} \\
\texttt{╟─┐◌┌─┼─╢ 6} \\
\texttt{║◌│◌│◌│◌║ 7} \\
\texttt{╚═╧═╧═╧═╝} \\
\\ 
\\ 
\texttt{Player 2:} \\
\texttt{ 1234567} \\
\texttt{╔══╤════╗} \\
\texttt{║T◌│D◌◌A║ 1} \\
\texttt{╟──┤◌◌┌─╢ 2} \\
\texttt{║◌◌├──┤◌║ 3} \\
\texttt{╟──┘◌◌├─╢ 4} \\
\texttt{║OP◌◌◌│S║ 5} \\
\texttt{╟─┐◌┌─┼─╢ 6} \\
\texttt{║◌│◌│M│◌║ 7} \\
\texttt{╚═╧═╧═╧═╝} \\
\\ 
\texttt{```} \\
            }
        }
    }
    & & \\ \\

    \theutterance \stepcounter{utterance}  
    & & & \multicolumn{2}{p{0.3\linewidth}}{
        \cellcolor[rgb]{0.9,0.9,0.9}{
            \makecell[{{p{\linewidth}}}]{
                \texttt{\tiny{[GM$|$GM]}}
                \texttt{End states:} \\
\texttt{```} \\
\texttt{Player 1:} \\
\texttt{ 1234567} \\
\texttt{╔══╤════╗} \\
\texttt{║◌◌│◌MA◌║ 1} \\
\texttt{╟──┤◌◌┌─╢ 2} \\
\texttt{║◌◌├──┤T║ 3} \\
\texttt{╟──┘◌◌├─╢ 4} \\
\texttt{║D◌◌◌◌│S║ 5} \\
\texttt{╟─┐◌┌─┼─╢ 6} \\
\texttt{║O│◌│P│◌║ 7} \\
\texttt{╚═╧═╧═╧═╝} \\
\\ 
\\ 
\texttt{Player 2:} \\
\texttt{ 1234567} \\
\texttt{╔══╤════╗} \\
\texttt{║◌◌│◌MAS║ 1} \\
\texttt{╟──┤◌◌┌─╢ 2} \\
\texttt{║◌◌├──┤T║ 3} \\
\texttt{╟──┘◌◌├─╢ 4} \\
\texttt{║D◌◌◌◌│◌║ 5} \\
\texttt{╟─┐◌┌─┼─╢ 6} \\
\texttt{║O│◌│P│◌║ 7} \\
\texttt{╚═╧═╧═╧═╝} \\
\\ 
\texttt{```} \\
            }
        }
    }
    & & \\ \\

    \theutterance \stepcounter{utterance}  
    & & & \multicolumn{2}{p{0.3\linewidth}}{
        \cellcolor[rgb]{0.9,0.9,0.9}{
            \makecell[{{p{\linewidth}}}]{
                \texttt{\tiny{[GM$|$GM]}}
                \texttt{* success: True} \\
\texttt{* lose: False} \\
\texttt{* aborted: False} \\
\texttt{{-}{-}{-}{-}{-}{-}{-}} \\
\texttt{* Shifts: 16.00} \\
\texttt{* Max Shifts: 12.00} \\
\texttt{* Min Shifts: 6.00} \\
\texttt{* End Distance Sum: 4.00} \\
\texttt{* Init Distance Sum: 16.49} \\
\texttt{* Expected Distance Sum: 29.33} \\
\texttt{* Penalties: 9.00} \\
\texttt{* Max Penalties: 16.00} \\
\texttt{* Rounds: 18.00} \\
\texttt{* Max Rounds: 28.00} \\
\texttt{* Object Count: 7.00} \\
            }
        }
    }
    & & \\ \\

\end{supertabular}
}

\end{document}
